% ---------------------------------------------------------------------
% The NeurIPS 2026 style file is embedded below so this single file
% compiles anywhere (Overleaf included) with no separate upload.
% DO NOT DELETE THIS BLOCK: it writes neurips_2026.sty at compile time.
% ---------------------------------------------------------------------
\begin{filecontents*}[overwrite]{neurips_2026.sty}
% partial rewrite of the LaTeX2e package for submissions to the
% Conference on Neural Information Processing Systems (NeurIPS):
% Roman Garnett (garnett@wustl.edu) and the many authors of
% nips15submit_e.sty, including MK and drstrip@sandia
% last revision: January 2026

\NeedsTeXFormat{LaTeX2e}
\ProvidesPackage{neurips_2026}[2026-01-29 NeurIPS 2026 submission/camera-ready style file]

\newif\if@neuripsfinal\@neuripsfinalfalse
\DeclareOption{final}{
  \@neuripsfinaltrue
  \@anonymousfalse
}

\newif\if@natbib\@natbibtrue
\DeclareOption{nonatbib}{
  \@natbibfalse
}

\newif\if@preprint\@preprintfalse
\DeclareOption{preprint}{
  \@preprinttrue
  \@anonymousfalse
}

\newif\if@main\@maintrue
\DeclareOption{main}{
  \@maintrue
  \newcommand{\@trackname}{\@neuripsordinal\ Conference on Neural Information Processing Systems (NeurIPS \@neuripsyear).}
}
\newif\if@position\@positionfalse
\DeclareOption{position}{
  \@positiontrue
  \newcommand{\@trackname}{\@neuripsordinal\ Conference on Neural Information Processing Systems (NeurIPS \@neuripsyear). Position Paper Track.}
}
\newif\if@eandd\@eanddfalse
\DeclareOption{eandd}{
  \@eanddtrue
\if@neuripsfinal\@anonymousfalse\else\if@preprint\@anonymousfalse\else\@anonymoustrue\fi\fi
  \newcommand{\@trackname}{\@neuripsordinal\ Conference on Neural Information Processing Systems (NeurIPS \@neuripsyear). Track on Evaluations and Datasets.} 
}
\newif\if@creativeai\@creativeaifalse
\DeclareOption{creativeai}{
  \@creativeaitrue
  \@anonymousfalse
  \newcommand{\@trackname}{\@neuripsordinal\ Conference on Neural Information Processing Systems (NeurIPS \@neuripsyear). Creative AI Track.}
}
\newif\if@education\@educationfalse
\DeclareOption{education}{
  \@educationtrue
  \@anonymousfalse
  \newcommand{\@trackname}{\@neuripsordinal\ Conference on Neural Information Processing Systems (NeurIPS \@neuripsyear). Education Track.}
}

\newif\if@anonymous\@anonymoustrue

\newcommand{\@workshoptitle}{}
\newcommand{\workshoptitle}[1]{\renewcommand{\@workshoptitle}{#1}}

\newif\if@workshop\@workshopfalse
\DeclareOption{sglblindworkshop}{
  \@workshoptrue
  \@anonymousfalse
  \newcommand{\@trackname}{\@neuripsordinal\ Conference on Neural Information Processing Systems (NeurIPS \@neuripsyear). Workshop: \@workshoptitle.}
}
\DeclareOption{dblblindworkshop}{
  \@workshoptrue
  \newcommand{\@trackname}{\@neuripsordinal\ Conference on Neural Information Processing Systems (NeurIPS \@neuripsyear). Workshop: \@workshoptitle.}
}
\DeclareOption{nonanonymous}{
  \@anonymousfalse
}

\ProcessOptions\relax

\renewcommand{\rmdefault}{ptm}
\renewcommand{\sfdefault}{phv}

\newcommand{\@neuripsordinal}{40th}
\newcommand{\@neuripsyear}{2026}
\newcommand{\@neuripslocation}{Sydney}

\usepackage{environ}
\newcommand{\acksection}{\section*{Acknowledgments and Disclosure of Funding}}
\NewEnviron{ack}{%
  \acksection
  \BODY
}

\if@natbib
  \RequirePackage{natbib}
\fi

\usepackage[verbose=true,letterpaper]{geometry}
\AtBeginDocument{
  \newgeometry{
    textheight=9in,
    textwidth=5.5in,
    top=1in,
    headheight=12pt,
    headsep=25pt,
    footskip=30pt
  }
  \@ifpackageloaded{fullpage}
    {\PackageWarning{neurips_2026}{fullpage package not allowed! Overwriting formatting.}}
    {}
}

\widowpenalty=10000
\clubpenalty=10000
\flushbottom
\sloppy

\renewcommand{\normalsize}{%
  \@setfontsize\normalsize\@xpt\@xipt
  \abovedisplayskip      7\p@ \@plus 2\p@ \@minus 5\p@
  \abovedisplayshortskip \z@ \@plus 3\p@
  \belowdisplayskip      \abovedisplayskip
  \belowdisplayshortskip 4\p@ \@plus 3\p@ \@minus 3\p@
}
\normalsize
\renewcommand{\small}{%
  \@setfontsize\small\@ixpt\@xpt
  \abovedisplayskip      6\p@ \@plus 1.5\p@ \@minus 4\p@
  \abovedisplayshortskip \z@  \@plus 2\p@
  \belowdisplayskip      \abovedisplayskip
  \belowdisplayshortskip 3\p@ \@plus 2\p@   \@minus 2\p@
}
\renewcommand{\footnotesize}{\@setfontsize\footnotesize\@ixpt\@xpt}
\renewcommand{\scriptsize}{\@setfontsize\scriptsize\@viipt\@viiipt}
\renewcommand{\tiny}{\@setfontsize\tiny\@vipt\@viipt}
\renewcommand{\large}{\@setfontsize\large\@xiipt{14}}
\renewcommand{\Large}{\@setfontsize\Large\@xivpt{16}}
\renewcommand{\LARGE}{\@setfontsize\LARGE\@xviipt{20}}
\renewcommand{\huge}{\@setfontsize\huge\@xxpt{23}}
\renewcommand{\Huge}{\@setfontsize\Huge\@xxvpt{28}}

\renewcommand{\tiny}{\fontsize{6pt}{7pt}\selectfont}
\renewcommand{\scriptsize}{\fontsize{7pt}{8pt}\selectfont}
\renewcommand{\footnotesize}{\fontsize{8pt}{9.5pt}\selectfont}

\providecommand{\section}{}
\renewcommand{\section}{%
  \@startsection{section}{1}{\z@}%
                {-2.0ex \@plus -0.5ex \@minus -0.2ex}%
                { 1.5ex \@plus  0.3ex \@minus  0.2ex}%
                {\large\bf\raggedright}%
}
\providecommand{\subsection}{}
\renewcommand{\subsection}{%
  \@startsection{subsection}{2}{\z@}%
                {-1.8ex \@plus -0.5ex \@minus -0.2ex}%
                { 0.8ex \@plus  0.2ex}%
                {\normalsize\bf\raggedright}%
}
\providecommand{\subsubsection}{}
\renewcommand{\subsubsection}{%
  \@startsection{subsubsection}{3}{\z@}%
                {-1.5ex \@plus -0.5ex \@minus -0.2ex}%
                { 0.5ex \@plus  0.2ex}%
                {\normalsize\bf\raggedright}%
}
\providecommand{\paragraph}{}
\renewcommand{\paragraph}{%
  \@startsection{paragraph}{4}{\z@}%
                {1.5ex \@plus 0.5ex \@minus 0.2ex}%
                {-1em}%
                {\normalsize\bf}%
}
\providecommand{\subparagraph}{}
\renewcommand{\subparagraph}{%
  \@startsection{subparagraph}{5}{\z@}%
                {1.5ex \@plus 0.5ex \@minus 0.2ex}%
                {-1em}%
                {\normalsize\bf}%
}
\providecommand{\subsubsubsection}{}
\renewcommand{\subsubsubsection}{%
  \vskip5pt{\noindent\normalsize\rm\raggedright}%
}

\renewcommand{\topfraction      }{0.85}
\renewcommand{\bottomfraction   }{0.4}
\renewcommand{\textfraction     }{0.1}
\renewcommand{\floatpagefraction}{0.7}

\newlength{\@neuripsabovecaptionskip}\setlength{\@neuripsabovecaptionskip}{7\p@}
\newlength{\@neuripsbelowcaptionskip}\setlength{\@neuripsbelowcaptionskip}{\z@}

\setlength{\abovecaptionskip}{\@neuripsabovecaptionskip}
\setlength{\belowcaptionskip}{\@neuripsbelowcaptionskip}

\renewenvironment{table}
  {\setlength{\abovecaptionskip}{\@neuripsbelowcaptionskip}%
   \setlength{\belowcaptionskip}{\@neuripsabovecaptionskip}%
   \@float{table}}
  {\end@float}

\setlength{\footnotesep }{6.65\p@}
\setlength{\skip\footins}{9\p@ \@plus 4\p@ \@minus 2\p@}
\renewcommand{\footnoterule}{\kern-3\p@ \hrule width 12pc \kern 2.6\p@}
\setcounter{footnote}{0}

\setlength{\parindent}{\z@}
\setlength{\parskip  }{5.5\p@}

\setlength{\topsep       }{4\p@ \@plus 1\p@   \@minus 2\p@}
\setlength{\partopsep    }{1\p@ \@plus 0.5\p@ \@minus 0.5\p@}
\setlength{\itemsep      }{2\p@ \@plus 1\p@   \@minus 0.5\p@}
\setlength{\parsep       }{2\p@ \@plus 1\p@   \@minus 0.5\p@}
\setlength{\leftmargin   }{3pc}
\setlength{\leftmargini  }{\leftmargin}
\setlength{\leftmarginii }{2em}
\setlength{\leftmarginiii}{1.5em}
\setlength{\leftmarginiv }{1.0em}
\setlength{\leftmarginv  }{0.5em}
\def\@listi  {\leftmargin\leftmargini}
\def\@listii {\leftmargin\leftmarginii
              \labelwidth\leftmarginii
              \advance\labelwidth-\labelsep
              \topsep  2\p@ \@plus 1\p@    \@minus 0.5\p@
              \parsep  1\p@ \@plus 0.5\p@ \@minus 0.5\p@
              \itemsep \parsep}
\def\@listiii{\leftmargin\leftmarginiii
              \labelwidth\leftmarginiii
              \advance\labelwidth-\labelsep
              \topsep    1\p@ \@plus 0.5\p@ \@minus 0.5\p@
              \parsep    \z@
              \partopsep 0.5\p@ \@plus 0\p@ \@minus 0.5\p@
              \itemsep \topsep}
\def\@listiv {\leftmargin\leftmarginiv
              \labelwidth\leftmarginiv
              \advance\labelwidth-\labelsep}
\def\@listv  {\leftmargin\leftmarginv
              \labelwidth\leftmarginv
              \advance\labelwidth-\labelsep}
\def\@listvi {\leftmargin\leftmarginvi
              \labelwidth\leftmarginvi
              \advance\labelwidth-\labelsep}

\providecommand{\maketitle}{}
\renewcommand{\maketitle}{%
  \par
  \begingroup
    \renewcommand{\thefootnote}{\fnsymbol{footnote}}
    \renewcommand{\@makefnmark}{\hbox to \z@{$^{\@thefnmark}$\hss}}
    \long\def\@makefntext##1{%
      \parindent 1em\noindent
      \hbox to 1.8em{\hss $\m@th ^{\@thefnmark}$}##1
    }
    \thispagestyle{empty}
    \@maketitle
    \@thanks
    \@notice
  \endgroup
  \let\maketitle\relax
  \let\thanks\relax
}

\newcommand{\@toptitlebar}{
  \hrule height 4\p@
  \vskip 0.25in
  \vskip -\parskip%
}
\newcommand{\@bottomtitlebar}{
  \vskip 0.29in
  \vskip -\parskip
  \hrule height 1\p@
  \vskip 0.09in%
}

\providecommand{\@maketitle}{}
\renewcommand{\@maketitle}{%
  \vbox{%
    \hsize\textwidth
    \linewidth\hsize
    \vskip 0.1in
    \@toptitlebar
    \centering
    {\LARGE\bf \@title\par}
    \@bottomtitlebar
    \if@anonymous
      \begin{tabular}[t]{c}\bf\rule{\z@}{24\p@}
        Anonymous Author(s) \\
        Affiliation \\
        Address \\
        \texttt{email} \\
      \end{tabular}%
    \else
      \def\And{%
        \end{tabular}\hfil\linebreak[0]\hfil%
        \begin{tabular}[t]{c}\bf\rule{\z@}{24\p@}\ignorespaces%
      }
      \def\AND{%
        \end{tabular}\hfil\linebreak[4]\hfil%
        \begin{tabular}[t]{c}\bf\rule{\z@}{24\p@}\ignorespaces%
      }
      \begin{tabular}[t]{c}\bf\rule{\z@}{24\p@}\@author\end{tabular}%
    \fi
    \vskip 0.3in \@minus 0.1in
  }
}

\newcommand{\ftype@noticebox}{8}
\newcommand{\@notice}{%
  \enlargethispage{2\baselineskip}%
  \@float{noticebox}[b]%
    \footnotesize\@noticestring%
  \end@float%
}

\renewenvironment{abstract}%
{%
  \vskip 0.075in%
  \centerline%
  {\large\bf Abstract}%
  \vspace{0.5ex}%
  \begin{quote}%
}
{
  \par%
  \end{quote}%
  \vskip 1ex%
}

\newcommand{\answerYes}[1][]{\textcolor{blue}{[Yes]#1}}
\newcommand{\answerNo}[1][]{\textcolor{orange}{[No]#1}}
\newcommand{\answerNA}[1][]{\textcolor{gray}{[N/A]#1}}
\newcommand{\answerTODO}[1][]{\textcolor{red}{\bf [TODO]}}
\newcommand{\justificationTODO}[1][]{\textcolor{red}{\bf [TODO]}}

\if@preprint
  \newcommand{\@noticestring}{%
    Preprint.%
  }
\else
  \if@neuripsfinal
    \newcommand{\@noticestring}{
      \@trackname
    }
  \else
    \newcommand{\@noticestring}{%
      Submitted to \@neuripsordinal\/ Conference on Neural Information Processing Systems (NeurIPS \@neuripsyear)\ifx\@workshoptitle\@empty\else. Workshop: \@workshoptitle\fi. Do not distribute.%
    }
    \NewEnviron{hide}{}
    \let\ack\hide
    \let\endack\endhide
    \RequirePackage{lineno}
    \linenumbers
    \AtBeginDocument{%
      \@ifpackageloaded{amsmath}{%
        \newcommand*\patchAmsMathEnvironmentForLineno[1]{%
          \expandafter\let\csname old#1\expandafter\endcsname\csname #1\endcsname
          \expandafter\let\csname oldend#1\expandafter\endcsname\csname end#1\endcsname
          \renewenvironment{#1}%
                          {\linenomath\csname old#1\endcsname}%
                          {\csname oldend#1\endcsname\endlinenomath}%
        }%
        \newcommand*\patchBothAmsMathEnvironmentsForLineno[1]{%
          \patchAmsMathEnvironmentForLineno{#1}%
          \patchAmsMathEnvironmentForLineno{#1*}%
        }%
        \patchBothAmsMathEnvironmentsForLineno{equation}%
        \patchBothAmsMathEnvironmentsForLineno{align}%
        \patchBothAmsMathEnvironmentsForLineno{flalign}%
        \patchBothAmsMathEnvironmentsForLineno{alignat}%
        \patchBothAmsMathEnvironmentsForLineno{gather}%
        \patchBothAmsMathEnvironmentsForLineno{multline}%
      }
      {}
    }
  \fi
\fi

\endinput
\end{filecontents*}
% ODI 2026 double-blind workshop revision (main2.tex, integration pass 2, 2026-09-05).
% Every number is either a value from main.tex or a value from a revision analysis whose independent check
% reports CONFIRMED or CONFIRMED_WITH_CAVEATS (I-02, I-02b, I-03, I-04, I-04b, I-05, I-05b, I-06, I-07, I-07b,
% I-08, I-09, I-10, I-12, I-16, I-21); those lines carry "% from analyses/I-xx_result.md; checked in I-xx_check.md".
% No test count is quoted (I-16_check.md: freeze-time run only). No pending or unsourced markers remain.
% Archived latency figures (the archived run) are labelled archived and are never attached to the clean configuration.
% Presentation polish: Tables 1-4 in the body at 8 pt, latency table in Appendix F, experiment dates restored by GPT polish check.
\documentclass{article}

\usepackage[dblblindworkshop]{neurips_2026}          % anonymous by default; [final] for camera-ready
\workshoptitle{On-Device Intelligence: Foundation Models under Real-World Constraints}

\usepackage[utf8]{inputenc}
\usepackage[T1]{fontenc}
\usepackage[hidelinks]{hyperref}
\usepackage{url}
\usepackage{booktabs}
\usepackage{amsmath}
\usepackage{amssymb}
\usepackage{graphicx}
\usepackage{microtype}
\usepackage{xcolor}
\usepackage{array}

\newcommand{\fn}[1]{\texttt{#1}}
\setlength{\textfloatsep}{2pt plus 1pt minus 1pt}
\setlength{\abovecaptionskip}{3pt}

\title{Fast Local AAC Board Construction:\\Diversity, Coverage and Runtime Trade-offs}
\author{Anonymous Author(s)}

\begin{document}
\maketitle

\begin{abstract}
Context-aware augmentative and alternative communication (AAC) requires timely suggestions while preserving user choice. We study a local architecture that separates speech and scene perception from board construction, replacing per-board language-model generation with fixed-vocabulary candidate filtering and an explicit selector that combines pinned core entries, answer reservations, a hedge cap and maximal marginal relevance (MMR). Across 120 synthetic scenarios with fixed teacher-generated candidate pools, MMR increases distinct communicative-function labels by 0.73 among six dynamic entries relative to relevance-only selection (95\% CI 0.58--0.88). % I-06_check.md
In a separate retrospective evaluation, learned and cosine candidate filters cover 61 and 60 of 71 annotated rows under the same selector, providing no established coverage advantage for the learned head. % I-02b_check.md
An OpenVINO serving configuration constructs boards in 4.50\,ms median on an Apple M4 CPU against 9.05\,ms for the PyTorch path, excluding perception, transport and display; the reduction comes from the embedding runtime, while head quantisation reduces head storage by 2.8$\times$. That configuration covers 62 of 71 rows against the teacher's 68, without establishing 90\% coverage retention. % I-10_check.md; I-04b_check.md
These results characterise component-level trade-offs in local AAC board construction; retrospective evaluation and proxy metrics limit generalisation, and benefit for AAC users remains untested.
\end{abstract}

% Related-work comparison checked against https://aclanthology.org/2021.acl-long.108.pdf sections 3.4--3.5
% and https://shawnshenjx.github.io/files/Kwickchat.pdf section 10 (two judges, 400 sentences).
\section{Introduction}
Context-aware augmentative and alternative communication (AAC) systems offer words or phrases, the \emph{entries} of a displayed board, that a user can select during an interaction. The board-construction stage must balance contextual relevance, a small number of available choices and timely updates. Language-model generation can supply flexible suggestions, but using generation for every board adds computation and gives the model substantial influence over which choices appear.

Context-aware AAC has drawn vocabulary from photographs by phrase clustering and frequency ranking \citep{fontana2021}, built context-sensitive tools \citep{kane2012}, expanded abbreviations \citep{cai2022}, generated sentences from keywords with a context-conditioned language model \citep{kwickchat2022} and predicted communication cards \citep{praact2024}; users report tensions between reduced typing and control over generated language \citep{valencia2023}. We differ in using transcript and scene context with explicit communicative-function constraints, and in selecting short fixed-vocabulary entries without run-time generation.

We investigate an architecture that separates contextual perception from constrained board construction. Speech and scene models run on a local laptop host; a lightweight board stage filters candidate entries and selects six dynamic suggestions alongside four pinned core entries. The wearable supplies audio and images and performs no inference. Candidate filters and inference runtimes can then change while explicit rules over the displayed choices are retained. Per-board generation with the local teacher took 7.05\,s per board (Appendix~\ref{app:additional}); the board stage studied here takes milliseconds.

Our question is how candidate selection and runtime choices affect measured diversity, coverage and post-perception latency. We examine it through a controlled selector comparison on fixed teacher-generated pools, retrospective comparisons of learned and non-learned candidate filters, and profiling of the serving pipeline, and contribute three empirical findings. First, MMR increases function-label diversity relative to relevance-only selection when candidate pools and board size are fixed (Section~\ref{sec:slate}). % from analyses/I-06_result.md; checked in I-06_check.md
Second, the learned candidate filter has no established coverage advantage over cosine filtering in the retrospective 95-entry comparison (Section~\ref{sec:distil}). % from analyses/I-02b_result.md; checked in I-02b_check.md
Third, the measured board-stage latency reduction comes from the embedding runtime, while head quantisation reduces storage without improving board-stage latency (Section~\ref{sec:latency}). % from analyses/I-10_result.md; checked in I-10_check.md
These findings characterise a local systems design; they do not establish improved communication outcomes for AAC users.

% Unchanged architecture/configuration constants: revision notes (issues I-08/I-11/I-17/I-23/I-27). Revision experiment values are sourced at their tables/paragraphs.
\section{Local architecture and inference boundary}\label{sec:system}
\begin{figure}[t]
\centering\footnotesize
\fbox{\begin{minipage}[c]{0.18\linewidth}\centering Wearable\\Audio + images\\No inference\end{minipage}}
$\longrightarrow$
\fbox{\begin{minipage}[c]{0.45\linewidth}\centering Local laptop host: speech + scene perception\\$\downarrow$\\\textbf{Board-construction boundary:} embed $\rightarrow$ score $\rightarrow$ select, \textbf{4.5\,ms median}\end{minipage}} % from analyses/I-10_result.md (clean configuration, OpenVINO embedder + INT8 head); checked in I-10_check.md
$\longrightarrow$
\fbox{\begin{minipage}[c]{0.16\linewidth}\centering Tablet\\4 core entries\\6 dynamic entries\end{minipage}}
\caption{The local network spans three devices. Reported board timings cover only the lower host stage on an Apple M4 laptop CPU (size-matched 95-entry configuration, Table~\ref{tab:latencymain}); perception, transport and display are outside the timed boundary.}
\label{fig:system}
\end{figure}
An ESP32-S3 wearable streams frames and audio over a WebSocket and performs no inference; every model runs on a laptop-class host (Apple M4 CPU in all measurements), which runs speech and scene perception (Whisper \citep{whisper}; Qwen2.5-VL-3B \citep{qwen25vl}) and adds personal context (Appendix~\ref{app:hw}). One function, \fn{rank\_vocabulary}, receives scene, transcript, personal context and a strategy and returns ranked entries; every strategy passes through it, and runtimes are chosen by configuration without changing firmware or the interface (Appendix~\ref{app:repro}). Buffers are ephemeral, personal context stays on the host, four core entries stay outside model control and an utterance requires a user selection; access control, transport security and bystander consent remain deployment requirements.

\section{Diversity-first board selection}\label{sec:slate}
\textbf{Policy.} A local Qwen2.5-3B-Instruct teacher generates candidates from a slot prompt covering nine communicative moves plus four topical phrases; a deterministic pool adds the four pinned core entries (Appendix~\ref{app:metrics}). Dynamic slots are filled by maximal marginal relevance (MMR) \citep{mmr1998} over \fn{all-MiniLM-L6-v2} embeddings \citep{minilm}, trading clipped cosine relevance to the context against similarity to entries already chosen ($\lambda=0.6$; answer-slot and tie-break rules in Appendix~\ref{app:train}); answer slots are reserved first, so MMR is a primitive the selector calls, not a whole-board optimisation.

\textbf{Protocol.} We use 120 synthetic scenarios: 24 each in classroom, clinic, airport, dinner and direct-question settings. We lead with dyn@6, the six entries the selector actually changes, keep @10 for the whole board and report the historical @7 window in the appendix. % from analyses/I-03_result.md; checked in I-03_check.md
Metrics (Appendix~\ref{app:metrics}): distinct nearest-exemplar function labels, near-duplicate pairs at cosine $\geq0.85$, proxy coverage of an acceptable move, and a lexical hit on the gold utterance. Table~\ref{tab:selection} is a same-pool comparison: one teacher pool per scenario was regenerated with the same model (greedy decoding, in a separate later run), and three arms fill the six dynamic slots from that pool with the same pinned entries and board size (Appendix~\ref{app:uncertainty}); the confounded comparison with direct LLM generation is in Table~\ref{tab:selectionfull}.

\begin{table}[t]
\centering\footnotesize\renewcommand{\arraystretch}{0.9}
\setlength{\tabcolsep}{5pt}\setlength{\belowcaptionskip}{3pt}
\caption{Same-pool selector comparison on 120 scenarios: same regenerated pool, pinned core and board size per scenario; arms differ only in filling the six dynamic slots. @7 rows and the direct-generation comparison: Table~\ref{tab:selectionfull}.}
\label{tab:selection}
\begin{tabular}{llrrrr}
\toprule
Window & Selection rule & Distinct fn.\ $\uparrow$ & Near-dup.\ $\downarrow$ & Coverage $\uparrow$ & Lexical hit\\
\midrule
% from analyses/I-06_result.md; checked in I-06_check.md (boards and scores reproduced byte-for-byte)
dyn@6 & Relevance only & 3.058 & 0.0306 & 0.958 & 0.608\\
dyn@6 & MMR ($\lambda=0.6$) & 3.783 & 0.0083 & 0.950 & 0.617\\
dyn@6 & Full shipped policy & 3.858 & 0.0078 & 0.958 & 0.575\\
@10 & Relevance only & 4.700 & 0.0104 & 0.992 & 0.608\\
@10 & MMR ($\lambda=0.6$) & 5.258 & 0.0030 & 1.000 & 0.617\\
@10 & Full shipped policy & 5.383 & 0.0026 & 1.000 & 0.575\\
\bottomrule
\end{tabular}
\end{table}

\textbf{Results.} With pools and board size fixed, plain MMR raises distinct functions over relevance-only selection by 0.73 $[0.58,0.88]$ at dyn@6, 0.56 $[0.43,0.68]$ at @10 and 0.43 $[0.32,0.55]$ at @7, and lowers near-duplicate and same-function pair rates (Appendix~\ref{app:uncertainty}). Because only the selection rule changes between arms, the effect is causal within a pool. The full shipped policy adds 0.08 $[-0.03,0.18]$, 0.13 $[0.04,0.21]$ and 0.03 $[-0.06,0.13]$ functions over plain MMR at the same windows. Coverage is near ceiling in every arm: relevance-only covers 115/120 scenarios, and MMR changes that by $-1/120$ $[-0.042,0.025]$ at dyn@6 and $+5/120$ $[0.008,0.083]$ at @7. % I-06_check.md
Mean clipped context relevance of the six chosen entries decreases from 0.354 under relevance-only selection to 0.317 under MMR and 0.282 under the full policy, a trade-off between embedding relevance and function-label diversity that the diversity metrics do not price; we do not interpret the decrease as a demonstrated loss of communication quality. % from analyses/I-06_result.md diagnostics; reproduced byte-identical in I-06_check.md
The complete-policy comparison is reported in Appendix~\ref{app:uncertainty} (Table~\ref{tab:selectionfull}(a)); its differences cannot be attributed to selection alone because the prompt, candidate pool, board length and pinning also differ. The same-pool experiment isolates selection effects on teacher-generated candidates; Section~\ref{sec:distil} separately evaluates fixed-vocabulary candidate filtering and the student selector. % I-03_check.md; I-06_check.md

\section{Distil the scores, retain the constraints}\label{sec:distil}
The teacher takes seconds per board (Appendix~\ref{app:additional}), so we replace generation with a student that scores a fixed vocabulary of $V$ entries: a two-layer head over the MiniLM embeddings of scene and transcript plus a question flag, trained to predict which entries the teacher selected (Appendix~\ref{app:train}). The head only chooses the candidate pool: its 18 highest-scoring entries go to the selector, whose MMR relevance term is context cosine.

\textbf{Corpus and vocabularies.} The teacher corpus has 600 boards (120 Section~\ref{sec:slate} contexts, 360 authored scenarios and 120 imported turns of undocumented provenance); deduplication and a 15\% category cap leave 559 rows, split 391/84/84. Evaluation boards never enter gradient updates, but 22 of the 84 evaluation rows share a template family or source conversation with a training row, and the partition was inspected during development (Section~\ref{sec:limits}). % from analyses/I-16_result.md and I-09_result.md; checked in I-16_check.md and I-09_check.md
We use three vocabularies with consistent names. \emph{Historical, 95 entries}: fitted on all 559 rows before splitting, so its label space is not independent of the evaluation contexts. \emph{Training-only, 45 entries}: vocabulary, canonical map, label priors and loss weights fitted on the 391 training rows only, same partition, hyperparameters and selector, unchanged frequency floor. % from analyses/I-05_result.md; checked in I-05_check.md
\emph{Size-matched, 95 entries}: the same training-only fit with the floor lowered and the count capped to the historical size; the cap was chosen after a floor-two diagnostic had been scored on the evaluation partition, so this is a retrospective control, not a validated configuration (Appendix~\ref{app:train}). % from analyses/I-05b_result.md; checked in I-05b_check.md
% Lexical representability 37.5 / 26.8 / 30.0 % of teacher non-core slots: Appendix D (analyses/I-16_result.md, I-05b_result.md)

\textbf{Selector.} A \emph{hedge} is one of 18 listed phrases that defer or express uncertainty (``I don't know'', ``maybe''); an \emph{open answer} is a span extracted from the context. From the 18-entry pool the student fills six dynamic slots with cosine MMR; on questions it first reserves two non-hedge candidates or up to two open answers, and at most one hedge is admitted during fill. The teacher has an answer quota but neither the student's non-hedge/open-answer reservation nor its hedge cap; on the historical checkpoint the selector places 0.19 listed hedges per seven-entry window against 1.07 for raw score-based selection and 0.43 for the teacher (Appendix~\ref{app:train}). % from analyses/I-08_result.md; checked in I-08_check.md
An appropriate ``I don't know'' is a legitimate move, so the cap is a policy choice; the selector also consults the evaluation taxonomy (Section~\ref{sec:limits}).

\begin{table}[t]
\centering\footnotesize\renewcommand{\arraystretch}{0.9}
\setlength{\tabcolsep}{5pt}\setlength{\belowcaptionskip}{3pt}
\caption{Teacher and students on the 84 evaluation-partition rows. \emph{Size-matched 95}: retrospective control with training-only label preprocessing; \emph{served}: the same checkpoint through the serving function (OpenVINO MiniLM, mixed-precision INT8 head calibrated on training contexts; other student rows use the direct PyTorch selector); \emph{training-only 45}: same fit, unchanged floor; \emph{historical}: vocabulary fitted before splitting. Coverage: 71 annotated rows, one measurement at @10 and dyn@6; agreement: teacher-entry recall. Identities: Appendix~\ref{app:repro}; intervals and @7 rows: Appendix~\ref{app:uncertainty}.}
\label{tab:student}
\begin{tabular}{llrrrrr}
\toprule
Window & Board & Dist.\ fn.\ $\uparrow$ & Near-dup.\ $\downarrow$ & Same-fn.\ $\downarrow$ & Dyn.\ cov.\ $\uparrow$ & Agree.\\
\midrule
% Teacher and historical rows: analyses/I-03_result.md, checked in I-03_check.md.
% Clean 95 rows (b72bb3ec7254): analyses/I-05b_result.md, checked in I-05b_check.md. Clean 45 rows (b691659cecf3): I-05_result.md, checked in I-05_check.md. Historical: d6ce2e56d50b.
dyn@6 & Teacher & 3.917 & 0.0127 & 0.1833 & 0.958 & ---\\
dyn@6 & Student, size-matched 95 & 4.869 & 0.0000 & 0.0857 & 0.859 & 0.079\\
dyn@6 & Student, size-matched 95, served & 4.917 & 0.0000 & 0.0825 & 0.873 & 0.077\\ % from analyses/I-04b_result.md; checked in I-04b_check.md
dyn@6 & Student, historical & 4.786 & 0.0000 & 0.0937 & 0.859 & 0.091\\
@10 & Teacher & 5.190 & 0.0042 & 0.1643 & 0.958 & ---\\
@10 & Student, size-matched 95 & 5.310 & 0.0000 & 0.1582 & 0.859 & 0.079\\
@10 & Student, size-matched 95, served & 5.417 & 0.0000 & 0.1529 & 0.873 & 0.077\\ % from analyses/I-04b_result.md (seam boards, I-04b_scores_v1.json); checked in I-04b_check.md
@10 & Student, training-only 45 & 5.036 & 0.0000 & 0.1714 & 0.761 & 0.071\\
@10 & Student, historical & 5.286 & 0.0000 & 0.1619 & 0.859 & 0.091\\
% @7 continuity rows (teacher 4.298/0.0006/0.1735/0.873; historical student 4.095/0.0000/0.1893/0.789/0.067) in Appendix E text.
\bottomrule
\end{tabular}
\end{table}

\textbf{Student results.} The size-matched 95-entry student covers fewer rows than the teacher and more function labels. Proxy acceptable-move coverage asks whether a board offers at least one acceptable move (71 annotated rows): the size-matched student covers 61/71 against the teacher's 68/71 (ratio 0.897; margin over 90\% of teacher $-0.003$, 95\% paired CI $[-0.100,0.090]$), the same count as the historical model, so the 90\% criterion remains inconclusive. The training-only 45-entry student covers 54/71 (ratio 0.794), and the paired size-matched minus training-only difference is $+0.099$ $[0.014,0.183]$; that contrast conflates vocabulary size, content and retraining and does not decompose leakage, and the @7 comparison with the historical model is in Appendix~\ref{app:uncertainty}. % from analyses/I-05_result.md and I-05b_result.md; checked in I-05_check.md and I-05b_check.md
The student's distinct-function point estimates exceed the teacher's at dyn@6 and @10 (margins over 90\% of teacher 1.344 $[1.121,1.556]$ and 0.638 $[0.385,0.889]$), the same-function-pair difference at @10 is inconclusive (0.006 $[-0.011,0.022]$), and pairs are fewer at dyn@6. Higher diversity accompanies different constraints and fewer hedges; it does not establish better meaning preservation. Exact-entry recall is low for every model (0.079 size-matched, 0.091 historical at @10); Table~\ref{tab:example} shows a matching function count beside a missed acceptable move. % from analyses/I-05b_result.md; checked in I-05b_check.md

% Worked example: mac/evaluation/phase4/teacher_data.jsonl id "dinner-4-1" (gathered row index 81, test split,
% category dinner, is_question true). Teacher board = the cached row["board"]; student board regenerated in revision
% with infer.student_board(model, vocab, scene, transcript, True, k=10, k_dyn=6, core) from the
% historical student.pt (config_hash d6ce2e56d50b, 95-entry vocab.json). Scores with metrics3 on that pair: teacher
% @7 distinct functions 4 / same-fn 0.190, student @7 4 / 0.286; coverage teacher True, student False at @7 and @10.
% Its three template siblings dinner-4-0/2/3 are in the training split (analyses/I-09_family_map.json).
\begin{table}[t]
\centering\footnotesize\renewcommand{\arraystretch}{0.9}\setlength{\belowcaptionskip}{3pt}
\caption{Worked example, evaluation row \fn{dinner-4-1} (three training siblings), historical student: the teacher covers an acceptable move, the student covers none, yet both score four distinct functions at @7.}
\label{tab:example}
\begin{tabular}{@{}p{0.115\linewidth}p{0.85\linewidth}@{}}
\toprule
Context & \emph{Scene:} A dinner table at home. dad is passing dishes around. \emph{Transcript:} dad: What do you want to do after dinner? \emph{Acceptable moves:} watch a movie; play a game; go to bed; read a book; go outside (the hedge move ``I don't know'' is excluded)\\
Teacher board & yes; no; I need help; wait; I want to cook; I'm not sure; I'd like to play video games; I need to rest; \textbf{I want to read a book}; I don't care\\
Student board & yes; no; I need help; wait; i'm not hungry; can i have some bread?; yes, please; i need to use the restroom.; i forgot; i need to check my schedule\\
\bottomrule
\end{tabular}
\end{table}

\textbf{Learned versus non-learned candidate filters.} Candidate filtering shows metric-dependent trade-offs under the same selector. Because the head only chooses the 18-entry pool, we replaced it under the identical selector with candidate filters without task-specific training: context cosine over the vocabulary, training-label frequency, the full vocabulary and a random pool (Appendix~\ref{app:uncertainty} reports all arms and intervals). On the size-matched 95-entry vocabulary the learned and cosine filters cover 61 and 60 of 71 annotated rows (difference 0.014, 95\% CI $-0.085$ to $0.113$). The learned filter yields 0.25 more distinct function labels among the six dynamic entries (4.869 against 4.619; uncorrected 95\% CI 0.012--0.476), while the full-board diversity difference is inconclusive ($+0.27$ for cosine, $[-0.01,0.55]$); cosine filtering produces 0.42 fewer listed hedges per board. Frequency covers 56/71 and a random pool 49/71; the random-minus-learned coverage difference is $-0.17$ $[-0.28,-0.06]$. % from analyses/I-02b_result.md; checked in I-02b_check.md (dyn@6 DF difference borderline and uncorrected among 45 intervals)
On the historical and training-only vocabularies the cosine filter raises full-board distinct functions ($+0.56$ and $+0.68$) and covers 65/71 against the learned 61/71 and 63/71 against 54/71, with frequency at 55/71 and 54/71. % from analyses/I-02_result.md; checked in I-02_check.md
These comparisons establish neither an overall advantage for learned filtering nor equivalence between the filters, and they do not allocate every quality outcome to individual components; for a local system the head is one compact candidate source among several. % from analyses/I-02_check.md and I-02b_check.md

\textbf{Export, quantisation and served-board quality.} Quantisation reduces head storage while preserving row-level proxy coverage at @10 and dyn@6 on this evaluation partition. Both training-fitted checkpoints pass OpenVINO FP32 parity; NNCF quantisation \citep{nncf} uses the 391 training contexts only and, by a rule chosen on the 84 validation contexts, excludes the final projection, giving a mixed-precision INT8 head: evaluation-context mean/minimum top-six overlap with FP32 is 0.950/0.833 for the size-matched head, which shrinks from 662 to 236\,KiB (2.8$\times$, excluding the embedder and perception models; Appendix~\ref{app:additional}). The historical head was calibrated on 360 contexts including 51 evaluation contexts, so its exported-path results are transductive. % from analyses/I-07_result.md and I-07b_result.md; checked in I-07_check.md and I-07b_check.md
Loaded through \fn{rank\_vocabulary} with OpenVINO MiniLM and the mixed-precision head, the size-matched control's boards cover 62/71 rows against the teacher's 68/71 at @10/dyn@6 (90\% retention margin $+0.011$ $[-0.085,0.096]$, inconclusive) (function counts and pair rates in Table~\ref{tab:student}). INT8 changes the entry set on 24 of 84 rows, but coverage in these windows matches the PyTorch serving path row for row; the one-row difference from the direct selector's 61/71 is not an improvement but a hedge match on a row whose question flag the serving heuristic misses (Appendix~\ref{app:uncertainty}). % I-04b_check.md; I-10_check.md

\section{Profile the board path, not just the head}\label{sec:latency}
% Archived head medians: docs/VERIFIED_NUMBERS.md micro-latency row; archived board medians: mac/intel/artifacts/e2e_mac.json (the archived run).
% Revision benchmark: analyses/I-10_result.md (I-10_bench_v3 + v4, 55,188 timed boards), checked in I-10_check.md; joined to quality by I-04b_check.md and I-10_check.md.
\begin{table}[t]
\centering\footnotesize\renewcommand{\arraystretch}{0.9}\setlength{\tabcolsep}{5pt}\setlength{\belowcaptionskip}{3pt}
\caption{Board-construction latency (ms), size-matched configuration, Apple M4, 5,040 calls per row; Jaccard against the PyTorch path (identical boards); detail in Table~\ref{tab:latency}.}
\label{tab:latencymain}
\begin{tabular}{llrrr}
\toprule
Embedder & Head & Median & IQR & Jaccard (identical)\\
\midrule
% from analyses/I-10_result.md (I-10_bench_v3 + v4); checked in I-10_check.md
PyTorch & PyTorch & 9.05 & 8.40--10.32 & 1.000 (84/84)\\
PyTorch & OpenVINO INT8 & 9.15 & 8.41--10.57 & 0.926 (60/84)\\
OpenVINO & OpenVINO INT8 & \textbf{4.50} & 3.98--5.15 & 0.926 (60/84)\\
\bottomrule
\end{tabular}
\end{table}
The latency reduction comes from the embedding runtime, not the head. We time the board-construction call through \fn{rank\_vocabulary} (embedding, head and selector) for the 84 evaluation-partition scenarios on an Apple M4 CPU, excluding perception, capture, transport and rendering: 5,040 timed calls per backend over six order-counterbalanced sessions of ten repeats, static embeddings warm, context keys evicted after each board, background load recorded rather than excluded (Appendix~\ref{app:hw}). For the size-matched configuration, whose saved first-round boards equal the quality boards, the OpenVINO embedder with the mixed-precision head takes 4.50\,ms median against 9.05\,ms for the PyTorch path and 9.15\,ms for the same head behind the PyTorch embedder (Table~\ref{tab:latencymain}); the paired per-scenario median ratio is 2.04 (IQR 1.87--2.28). p90 values are descriptive of this scenario mix, and the archived run's larger PyTorch-path tail is protocol-sensitive and not attributed to either head (Appendix~\ref{app:additional}). % from analyses/I-10_result.md and I-04_check.md; checked in I-10_check.md (mechanism-neutral wording required)
Head micro-benchmarks show a size saving without a head-time reduction (Appendix~\ref{app:hw}). The user-visible budget is dominated by the stages before this call: on the same host, voice-activity detection waits 600\,ms of trailing silence, Whisper takes 250--410\,ms per clip on either runtime, and scene description adds 0.93\,s when the change detector fires (Appendix~\ref{app:hw}), so the board stage is a small fraction of a response of about one second, or two when the scene model fires. We quote no teacher-to-student speed-up because the workloads and cache states differ (Appendix~\ref{app:additional}).

\section{Limitations and conclusion}\label{sec:limits}
\textbf{Limitations.} The corpus mixes synthetic and authored scenarios with imported turns of undocumented provenance; corpus-wide context deduplication precedes the split for every model, and the historical label space was fitted before it. % I-15/I-16; I-05
The evaluation partition was inspected during selector and criterion development, 22 of its 84 rows share a template family or source conversation with a training row, and neither training-only preprocessing nor the size-matched control undoes that exposure; generalisation to unseen template families is untested, and the same-pool comparison uses one greedy teacher realisation per scenario (Appendix~\ref{app:uncertainty}). % I-14; I-09
Function diversity is an internal proxy: labels are unvalidated nearest-exemplar assignments, the selector consults the same taxonomy, and proxy acceptable-move coverage ignores polarity, so a negation can match its affirmative. % I-13
Timings are laptop-hosted, post-perception and warm-cache under recorded background load; a head-free board path, matched-workload teacher--student timing, concurrent load, energy, peak memory, cold start and controlled target-device runs are unmeasured (Appendix~\ref{app:additional}). % I-26; I-27; I-02_check.md; I-10_check.md
There is no evaluation with AAC users, so proxy diversity does not establish expressiveness, preference or successful communication; a consented, participatory evaluation is the priority. % I-26

\textbf{Conclusion.} This study characterises the trade-offs of separating local AAC board construction from language-model generation. On fixed teacher-generated pools, MMR increases function-label diversity among the six dynamic suggestions; under the constrained selector, the retrospective 95-entry comparison does not establish a coverage advantage for learned over cosine candidate filtering; and changing the embedding runtime reduces median board construction from 9.05 to 4.50\,ms on an Apple M4 CPU, while head quantisation primarily reduces storage. % from analyses/I-06, I-02b and I-10 results; checked in the matching _check.md files
The results also expose the limits of diversity as a quality measure: a board can represent several function labels yet omit an acceptable response (Table~\ref{tab:example}). The serving configuration covers 62 of 71 rows against 68 for the teacher with 90\% retention unestablished, so the contribution is an empirical account rather than evidence of improved communication outcomes. For local AAC design, these findings motivate evaluating candidate filtering, selection rules and inference runtime separately, with diversity assessed alongside acceptable-response coverage. % from analyses/I-04b_result.md; checked in I-04b_check.md

\clearpage

\begin{thebibliography}{99}

\bibitem[Bai et al.(2025)]{qwen25vl}
S.~Bai, K.~Chen, X.~Liu, J.~Wang, W.~Ge, S.~Song, K.~Dang, P.~Wang, S.~Wang,
J.~Tang, H.~Zhong, Y.~Zhu, M.~Yang, Z.~Li, J.~Wan, P.~Wang, W.~Ding, Z.~Fu,
Y.~Xu, J.~Ye, X.~Zhang, T.~Xie, Z.~Cheng, H.~Zhang, Z.~Yang, H.~Xu, and
J.~Lin.
\newblock Qwen2.5-VL technical report.
\newblock \emph{arXiv preprint arXiv:2502.13923}, 2025.

\bibitem[Cai et al.(2022)]{cai2022}
S.~Cai, S.~Venugopalan, K.~Tomanek, A.~Narayanan, M.~R. Morris, and M.~P.
Brenner.
\newblock Context-aware abbreviation expansion using large language models.
\newblock In \emph{NAACL-HLT}, pages 1261--1275, 2022.

\bibitem[Carbonell and Goldstein(1998)]{mmr1998}
J.~Carbonell and J.~Goldstein.
\newblock The use of MMR, diversity-based reranking for reordering documents and
producing summaries.
\newblock In \emph{SIGIR}, pages 335--336, 1998.

\bibitem[Fontana de Vargas and Moffatt(2021)]{fontana2021}
M.~Fontana de Vargas and K.~Moffatt.
\newblock Automated generation of storytelling vocabulary from photographs for
use in AAC.
\newblock In \emph{ACL-IJCNLP}, pages 1353--1364, 2021.

\bibitem[Kane et al.(2012)]{kane2012}
S.~K. Kane, B.~Linam-Church, K.~Althoff, and D.~McCall.
\newblock What we talk about: Designing a context-aware communication tool for
people with aphasia.
\newblock In \emph{ASSETS}, pages 49--56, 2012.

\bibitem[Kozlov et al.(2020)]{nncf}
A.~Kozlov, I.~Lazarevich, V.~Shamporov, N.~Lyalyushkin, and Y.~Gorbachev.
\newblock Neural network compression framework for fast model inference.
\newblock \emph{arXiv preprint arXiv:2002.08679}, 2020.

\bibitem[Pereira et al.(2024)]{praact2024}
Jayr Alencar Pereira, Jaylton Alencar Pereira, Cleber Zanchettin, and Robson do Nascimento Fidalgo.
\newblock PrAACT: Predictive augmentative and alternative communication with
transformers.
\newblock \emph{Expert Systems with Applications}, 240:122417, 2024.

\bibitem[Radford et al.(2023)]{whisper}
A.~Radford, J.~W. Kim, T.~Xu, G.~Brockman, C.~McLeavey, and I.~Sutskever.
\newblock Robust speech recognition via large-scale weak supervision.
\newblock In \emph{ICML}, volume 202 of \emph{PMLR}, pages 28492--28518, 2023.

\bibitem[Shen et al.(2022)]{kwickchat2022}
J.~Shen, B.~Yang, J.~J. Dudley, and P.~O. Kristensson.
\newblock KWickChat: A multi-turn dialogue system for AAC using context-aware
sentence generation by bag-of-keywords.
\newblock In \emph{IUI}, pages 853--867, 2022.

\bibitem[Valencia et al.(2023)]{valencia2023}
S.~Valencia, R.~Cave, K.~Kallarackal, K.~Seaver, M.~Terry, and S.~K. Kane.
\newblock ``The less I type, the better'': How AI language models can enhance
or impede communication for AAC users.
\newblock In \emph{CHI}, 2023.

\bibitem[Wang et al.(2020)]{minilm}
W.~Wang, F.~Wei, L.~Dong, H.~Bao, N.~Yang, and M.~Zhou.
\newblock MiniLM: Deep self-attention distillation for task-agnostic compression
of pre-trained transformers.
\newblock In \emph{NeurIPS}, volume~33, 2020.

\end{thebibliography}
\clearpage
\appendix
% =====================================================================
\section{Metric definitions}\label{app:metrics}

Full@10 scores up to ten entries including core; dyn@6 removes core entries case-insensitively before retaining at most six in each arm. Full@7 retains the historical first-seven window. Dynamic coverage in a full window removes core entries after truncation. % from analyses/I-03_result.md; checked in I-03_check.md
The pinned set (``yes'', ``no'', ``I need help'', ``wait'') is classified as comment, refuse/redirect, request-help and comment, so boards with these pins start with exactly three distinct functions and one same-function pair: a same-function-pair floor of $1/21$ at seven entries, $1/45$ at ten, and zero when only the six dynamic entries are scored.

\begin{description}
\item[Distinct functions] Number of the nine AAC communicative functions
(question-back, specific-answer, I-don't-know, request-help, comment,
refuse/redirect, emotion, request-clarification, social) represented among the
evaluated window. Each entry receives the label of its single nearest authored exemplar (maximum cosine, not nearest centroid) in the 384-d MiniLM embedding space. Generation tags instead use request, reject, comment, question, share-information, express-feeling, social, direct-attention and clarify; the evaluation labels are assigned independently of these tags.
\item[Near-duplicate pair rate] Fraction of entry pairs in the evaluated window whose embedding cosine similarity is at least 0.85.
\item[Same-function-pair rate] Fraction of entry pairs in the evaluated window that
share a function label. The distillation criterion is ``no worse than the teacher''; lower is not always better, because several concrete answers to one question share a label.
\item[Proxy acceptable-move coverage (``coverage'')] Fraction of scenarios for which at least one of the
scenario's authored \emph{acceptable moves} matches a board entry under the proxy rule.
The all-entry score uses the stated window. ``Dynamic'' removes core entries from that window before scoring and removes hedge moves from the acceptable-move set; hedge \emph{predictions} are not removed and can still receive lexical or cosine credit. Matching uses word-boundary lexical overlap or embedding cosine at least 0.50 and ignores polarity. Table~\ref{tab:selection} coverage keeps hedge targets. Dynamic coverage in Table~\ref{tab:student} is measured on the 71 test rows with authored acceptable moves (20 Section~\ref{sec:slate} rows and 51 authored scenarios; the 13 imported turns have none).
\item[Lexical hit (diagnostic)] Whether the scenario's single authored gold utterance is lexically matched by any entry of the full generated board. It is a lexical-hit test, not exact-entry matching.
\item[Board variety] Number of distinct unordered entry sets in the stated window over the evaluation partition divided by its size; a degeneracy guard. At @10 it is 0.976 for the size-matched 95-entry student, 0.964 for the training-only 45-entry student and 1.000 for the historical student. % from analyses/I-05b_result.md, I-05_result.md, I-03_result.md
\item[Dynamic word agreement] Directional exact-entry recall after lower-casing. At @10 and dyn@6 it compares the teacher's six non-core entries with the student's six. At @7 it retains the legacy asymmetric rule: teacher non-core entries within the first seven versus all six student dynamic entries. % from analyses/I-03_result.md; checked in I-03_check.md
\item[Diversity retention] Student distinct functions divided by the teacher's, in the same window.
% Hedge-list correction: analyses/I-08_check.md.
\item[Hedge] One of 18 listed phrases expressing uncertainty or deferral; the same list serves the selector's cap and the metrics. Near-hedges outside the list (``I need to think about it'') are neither capped nor counted.
\item[INT8 top-$k$ overlap] For each context, $|T_k^{\mathrm{INT8}}
\cap T_k^{\mathrm{FP32}}|/k$ with $k{=}6$; mean and minimum over the stated context set (the 360 calibration contexts for the historical head; the 84 test contexts for the training-only 45- and 95-entry heads).
\item[Board Jaccard] $|A\cap B|/|A\cup B|$ for a backend board $A$ and the PyTorch board $B$ on the same scenario, averaged over the 84 held-out rows.
\end{description}

\section{Reproduction}\label{app:repro}

% Environment-variable prefix anonymised for review; restore at camera-ready.
The commands below describe the historical workflow and explicit-window re-scoring, run from the host server directory with perception disabled (\fn{SYSTEM\_PERCEPTION=0}); none calls an LLM, VLM or Whisper unless noted. Running them requires the project code, models and evaluation data, which are planned for release on acceptance; the \fn{SYSTEM} prefix is an anonymised placeholder for the project prefix. Default-window scripts overwrite historical outputs; use a fresh output path for each explicit-window re-score. The MiniLM bundle must be exported and validated before OpenVINO benchmarking; a missing bundle is rejected by the benchmark, and the revision quality scripts assert generation backends before scoring in a separate PyTorch process. % from analyses/I-12_check.md
The revision analyses (training-only preprocessing and retraining, the size-matched control, the candidate-pool ablation, the same-pool selector ablation, train-only requantisation and the constant recounts) are stand-alone scripts that import the project modules and write hashed manifests; they will be released with the code.
\begin{footnotesize}
\begin{verbatim}
python3 -m pytest -q                                  # record the dated outcome
python3 evaluation/phase3/selftest.py --from-boards   # Table 5(a) (re-scores cached boards)
# Repeat with --metric-k 7 for continuity, or add --dynamic-only for dyn@6 (k=6).
python3 evaluation/phase3/selftest.py --from-boards \
  --metric-k 10 --output /tmp/slate10.json
python3 evaluation/phase4/evaluate.py --teacher-samples 0 \
  --metric-k 10 --output /tmp/student10.json
python3 evaluation/phase4/evaluate.py --teacher-samples 8  # teacher latency (LLM run)
python3 intel/convert.py                              # student IR + MiniLM export
python3 intel/convert.py --skip-student --validate-minilm-only # existing bundle check
python3 intel/quantize.py                             # NNCF INT8 + overlap check
python3 intel/bench.py                                # head micro-latency, benchmark_app
python3 intel/bench_e2e.py --n 84                     # latency table (post-perception board)
python3 intel/bench_whisper.py --n 10                 # Whisper comparison
\end{verbatim}
\end{footnotesize}
% No test count is asserted: I-16_check.md requires one dated run at repository freeze (the 2026-09-05 21:44 EDT run, 241 passed / 1 media-scan failure, is retained in revision/pytest_freeze_2026-09-05.log for the final build).

Artefact identities: historical configuration \fn{d6ce2e56d50b}, vocabulary \fn{865b8083da276824}; training-fitted 45-entry configuration \fn{b691659cecf3}, vocabulary \fn{30e9becb9a556cd8}; retrospective 95-entry configuration \fn{b72bb3ec7254}, vocabulary \fn{0f3bb094cd439111}. All use split fingerprint \fn{ab51f7fba4c9c39e}. The historical configuration ID is recorded provenance; the current parsed YAML has ID \fn{5c185c1fcb78}, and the differently shaped revision configuration objects do not imply that every setting changed. % I-05_check.md; I-05b_check.md; I-21_check.md

The shared SHA-256 archive (\fn{34e556f7}\ldots) covers the historical artefacts and the initial training-fitted 45-entry model, candidate controls and export. It predates the size-matched control, served-quality and revision-latency experiments, which have their own hashed records. It preserves working-tree bytes, not a source-control revision, and excludes external MiniLM and teacher weights. % I-21_check.md
Runtime selection uses \fn{SYSTEM\_STUDENT\_BACKEND=torch|openvino} and \fn{SYSTEM\_EMBED\_BACKEND=torch|openvino}; head device selection is \fn{SYSTEM\_OV\_DEVICE=CPU|GPU|NPU}. Vocabulary sidecars reject mismatched head/vocabulary pairs; split fingerprints detect corpus drift, not family overlap.

\section{Hardware and software fingerprint}\label{app:hw}


\begin{itemize}
\item Host for every number in Tables 1--5 and the abstract: Apple M4 (arm64),
Darwin 25.5.0. Exploratory Intel observations appear in Appendix~\ref{app:additional}, from one uncontrolled session on a Core Ultra 9 185H;
no controlled absolute Intel board latency is claimed.
\item OpenVINO 2026.2.1 (build 21919), ARM CPU plugin, devices [CPU],
\fn{INFERENCE\_PRECISION\_HINT=f32} (pinned because the ARM plugin's f16 default degraded Whisper decoding and shifted student logits); PyTorch 2.12.0; Python 3.13.6; NNCF 3.2.0.
\item Head micro-benchmarks: batch 1, input (1, 769) float32, 500 timed
iterations after 50 warm-up; \fn{benchmark\_app} latency hint, 5\,s each. Head-only medians on the Apple M4: 0.015\,ms (PyTorch), 0.025\,ms (OpenVINO FP32) and 0.031\,ms (INT8). % Archived head medians: docs/VERIFIED_NUMBERS.md micro-latency row (unchanged from main.tex)
\item Board benchmark (revision): 84 held-out scenarios; per backend six in-process sessions in a Latin square over backend order, run twice, each with the documented Python cache and seam-state reset, three untimed warm-up boards and ten timed repeats (5,040 calls), static embeddings pre-warmed and each board's context keys evicted after the call; plus one isolated single-backend process (840 calls) and one session of the archived recipe (three repeats, per-round eviction, 252 calls). Percentiles are over pooled samples; intervals resample the 84 scenarios. % from analyses/I-10_result.md; checked in I-10_check.md
\item Student artefacts on disk (\fn{.xml}+\fn{.bin}), historical 95-entry head: FP32 IR 662.1\,KiB (678,025 bytes), INT8 IR 235.5\,KiB (241,117 bytes). Training-only 45-entry head: FP32 IR 624.5\,KiB, INT8 IR 197.9\,KiB. Size-matched 95-entry head: FP32 IR 662.2\,KiB, INT8 IR 235.6\,KiB ($2.81\times$). % from analyses/I-07_result.md and I-07b_result.md; checked in I-07_check.md and I-07b_check.md (historical byte counts from I-07b_check.md)
\item Perception scheduling (Section~\ref{sec:system}): the wearable sends one VGA JPEG every 2\,s and 16\,kHz mono PCM; Whisper \fn{base.en} INT8 transcribes behind a voice activity detector, the scene model runs only when the change detector triggers, model work runs in worker threads so that scene generation stays off the audio path, and traffic stays within the local network.
\item Teacher: Qwen2.5-3B-Instruct (4-bit) via a local runtime; scene model
Qwen2.5-VL-3B-Instruct (4-bit); embedder \fn{all-MiniLM-L6-v2}; speech
Whisper \fn{base.en} INT8 via \fn{faster-whisper} (CTranslate2) or
\fn{openvino\_genai.WhisperPipeline}.
\item Wearable: ESP32-S3 (dual-core Xtensa LX7, 240\,MHz, 8\,MB PSRAM),
OV2640-class camera, SPH0645 I$^2$S microphone; VGA JPEG at 2000\,ms cadence,
16\,kHz PCM in 512-sample chunks. Separate FreeRTOS tasks isolate audio capture from camera work; a physical button toggles capture and an LED indicates activity. On the host, a 64$\times$64 greyscale change detector (archived gate timing 0.214\,ms per frame) gates scene inference. We quote no device cost, because no bill of materials has been costed.
\end{itemize}

\section{Training and selector configuration}\label{app:train}
The head is a two-layer network (769$\rightarrow$192$\rightarrow V$, dropout 0.3) over the concatenated 384-dimensional scene and transcript MiniLM embeddings and the question flag, trained with weighted binary cross-entropy, logit adjustment and an auxiliary listwise term (below). The MMR primitive of Section~\ref{sec:slate} is $s^*=\arg\max_{s\in C\setminus S}[\lambda\,\mathrm{rel}(s,c)-(1-\lambda)\max_{s'\in S}\mathrm{sim}(s,s')]$, where $C$ is the candidate pool, $S$ the set already chosen by the MMR pass, $c$ the context and $\mathrm{sim}$ cosine similarity; relevance is clipped cosine to the normalised context query, and an unseen function tag is preferred when scores tie within 0.02. On questions, tagged answers may occupy up to five dynamic slots and are chosen with $\lambda=0.35$; the remaining slots are filled with $\lambda=0.6$.

The training implementation uses Adam, learning rate $10^{-3}$, weight decay $10^{-3}$, batch size 32, seed 0, a maximum of 60 epochs and early stopping after 15 epochs without validation-loss improvement. The minimum-validation-loss checkpoint is retained. Historical vocabulary filtering uses minimum frequency three and a maximum of 200 entries before single-link embedding deduplication at cosine $\tau=0.80$ (23 merges); the training-only 45-entry model keeps those settings on the 391 training rows (58 raw entries, 45 after deduplication with 13 merges, validation selects epoch 6); the size-matched control uses minimum frequency two and a cap of 116 raw entries, the cap falling inside the frequency-two tier so that \fn{build\_vocab}'s existing alphabetical tie-break decides which frequency-two entries enter (95 after deduplication with 21 merges, 67 shared with the historical vocabulary, epoch 8). Corpus-level deduplication and the category cap still run on all rows. % from analyses/I-05_result.md and I-05b_result.md; checked
Lexical representability of the teacher's non-core entries (exact membership or one canonical-map lookup; 3600 corpus slots, 504 held-out slots): historical vocabulary 32.4\% (1166/3600; 23.4\% by exact match; 70/600 rows with none, 11.7\%) and 37.5\% (189/504) held out with 11/84 zero-hit rows; training-only 45-entry vocabulary 26.8\% (135/504) held out with 16/84 zero-hit rows; size-matched 95-entry control 30.0\% (151/504) held out. % from analyses/I-16_result.md and I-05b_result.md; checked in I-16_check.md and I-05b_check.md
Context deduplication uses cosine 0.95. The context query is a normalised weighted sum of the transcript, scene and personal-context embeddings (weights 0.7/0.2/0.1).

For label $j$, let $p_j$ be its empirical frequency, $z_j$ its logit and $y_j$ its binary target. The loss is weighted BCE on $z_j-\log p_j$, plus $0.3$ times the mean negative log-softmax probability of positive labels. Positive weights are $\mathrm{clip}(N_{\rm neg}/N_{\rm pos},1,20)$. For the historical model, vocabulary, priors and weights were computed over the pre-split corpus; for the training-fitted models they are fitted on training rows only, and poisoning every held-out board leaves all fitted objects unchanged. % from analyses/I-05_check.md
Gradient updates use training indices and checkpoint selection uses validation indices in every run.

The student scores a closed vocabulary, but its final board can include extracted open-answer spans, which may coincide with a vocabulary entry. Those spans receive up to two reserved slots on supported question types, in extractor order with a cosine-0.85 duplicate guard; reserved non-hedge candidates are chosen by query relevance with the same guard, and the hedge cap is one during fill and zero when an open answer is present. A small-pool backfill can relax the hedge cap to return six dynamic entries; it never fired on the historical checkpoint's held-out split. % from analyses/I-08_result.md; checked in I-08_check.md
Pinned core entries are assembled first. Reservation and MMR ordering are computed separately, and the MMR order is then filtered into the unreserved slots (Section~\ref{sec:slate}); the MMR relevance term is context cosine, so learned logits determine pool membership only.

On the 84 held-out rows of the historical checkpoint \fn{d6ce2e56d50b}, with the fixed 18-phrase hedge list, direct score-based selection (the four pinned entries plus the six highest-scoring entries) places 1.07 hedge entries in the seven-entry window and at least two on 7.1\% of boards (1.23 and 22.6\% over the whole board); the selector places 0.19 per window and 0.60 per board against the teacher's 0.43 and 0.98, and no student board carries more than one. % from analyses/I-08_result.md; checked in I-08_check.md
Canonicalisation maps six hedge variants to ``I don't know'' and one to ``I need more time'', leaving four vocabulary hedge entries. The same-checkpoint policy comparison does not measure a canonicalisation effect or isolate the cap. % from analyses/I-08_check.md (corrects I-08_result.md).

\section{Uncertainty and supplementary comparisons}\label{app:uncertainty}
\begin{table}[h]
\centering\footnotesize\renewcommand{\arraystretch}{0.9}
\setlength{\tabcolsep}{5pt}\setlength{\belowcaptionskip}{3pt}
\caption{Board selection on 120 scenarios, all windows. (a) Cached complete policies: direct generation returned eight entries per board, so @10 scores ten against eight and dyn@6 six against a mean of 5.94 (length-confounded); both arms satisfy the historical near-duplicate ceiling. (b) Same-pool ablation (pools regenerated in a separate later run, a different teacher realisation); @7 is the historical reporting window.}
\label{tab:selectionfull}
\begin{tabular}{llrrrr}
\toprule
Window & Policy & Distinct fn.\ $\uparrow$ & Near-dup.\ $\downarrow$ & Coverage $\uparrow$ & Lexical hit\\
\midrule
\multicolumn{6}{l}{\emph{(a) Complete policies, cached boards}}\\
% from analyses/I-03_result.md; checked in I-03_check.md (historical checkpoint).
@7 & Diversity-first & 4.683 & 0.0004 & 1.000 & 0.525\\
@7 & Direct LLM generation & 3.875 & 0.0274 & 0.775 & 0.333\\
@10 & Diversity-first & 5.425 & 0.0037 & 1.000 & 0.525\\
@10 & Direct LLM generation & 4.017 & 0.0292 & 0.775 & 0.333\\
dyn@6 & Diversity-first & 3.925 & 0.0106 & 0.975 & 0.525\\
dyn@6 & Direct LLM generation & 3.592 & 0.0292 & 0.650 & 0.333\\
\midrule
\multicolumn{6}{l}{\emph{(b) Same pool, same pinned core, same board size, historical @7 window}}\\
% from analyses/I-06_result.md; checked in I-06_check.md
@7 & Relevance only & 4.167 & 0.0095 & 0.958 & 0.608\\
@7 & MMR ($\lambda=0.6$) & 4.600 & 0.0012 & 1.000 & 0.617\\
@7 & Full shipped policy & 4.633 & 0.0004 & 0.983 & 0.575\\
\bottomrule
\end{tabular}
\end{table}

\begin{table}[h]
\centering\footnotesize\renewcommand{\arraystretch}{0.9}
\setlength{\tabcolsep}{5pt}\setlength{\belowcaptionskip}{3pt}
\caption{Candidate pools under the identical selector on the 84 test rows: distinct functions (rows covered, of 71). Learned: top-18 head logits; cosine: top-18 by context cosine; frequency: top-18 by training-label frequency (ties by NumPy quicksort over vocabulary positions); full: every entry. Historical is the pre-split vocabulary; training-only 45 and size-matched 95 are training-fitted.}
\label{tab:pools}
\begin{tabular}{llrrrr}
\toprule
Vocabulary & Window & Learned & Cosine & Frequency & Full vocabulary\\
\midrule
% Historical and clean 45: analyses/I-02_result.md, checked in I-02_check.md; clean 95: analyses/I-02b_result.md, checked in I-02b_check.md
% (pools, boards and per-row metrics reproduced independently in both checks); cells transposed from the previous pool-by-column layout
Historical 95 & @7 & 4.095 (56) & 4.667 (57) & 3.762 (44) & 4.690 (57)\\
Historical 95 & @10 & 5.286 (61) & 5.845 (65) & 4.881 (55) & 6.310 (64)\\
Historical 95 & dyn@6 & 4.786 (61) & 4.738 (65) & 4.679 (55) & 5.179 (64)\\
Training-only 45 & @7 & 3.833 (48) & 4.345 (53) & 3.512 (43) & 4.298 (53)\\
Training-only 45 & @10 & 5.036 (54) & 5.714 (63) & 4.333 (54) & 5.810 (61)\\
Training-only 45 & dyn@6 & 4.833 (54) & 4.917 (63) & 4.262 (54) & 5.024 (61)\\
Size-matched 95 & @7 & 4.179 (50) & 4.333 (48) & 4.012 (44) & 4.357 (47)\\
Size-matched 95 & @10 & 5.310 (61) & 5.583 (60) & 5.119 (56) & 5.762 (57)\\
Size-matched 95 & dyn@6 & 4.869 (61) & 4.619 (60) & 4.810 (56) & 4.738 (57)\\
\bottomrule
\end{tabular}
\end{table}

Quality intervals are 95\% percentile intervals from 10,000 paired row-bootstrap resamples (seed 0), conditional on the chosen corpus, policies and checkpoints. They omit family dependence, training-seed variation, configuration selection and multiplicity. Latency intervals instead use 2,000 scenario-bootstrap resamples, retaining all repeats of each sampled scenario (I-10). An interval containing zero does not establish equivalence.

\textbf{Student criteria.} Expressed as retention (student distinct functions over the teacher's), the size-matched 95-entry student scores 102.3\% at @10 and 124.3\% at dyn@6. These ratios do not isolate the effect of the selector constraints (Section~\ref{sec:distil}). A positive margin favours satisfying the criterion: for function diversity the margin is student minus $0.9$ times teacher; for same-function pairs it is teacher minus student; for coverage it is student minus $0.9$ times teacher.
\begin{center}\footnotesize
\setlength{\tabcolsep}{4pt}
\begin{tabular}{llrrr}
\toprule
Window & Student & Function diversity & Same-function pairs & Dynamic coverage\\
\midrule
% Clean 95: analyses/I-05b_result.md (checked in I-05b_check.md); clean 45: I-05_result.md (I-05_check.md); historical: I-03_result.md (I-03_check.md).
@10 & Size-matched 95 & 0.638 $[0.385,0.889]$ & 0.006 $[-0.011,0.022]$ & $-0.003$ $[-0.100,0.090]$\\
@10 & Training-only 45 & 0.364 $[0.131,0.592]$ & $-0.007$ $[-0.023,0.008]$ & $-0.101$ $[-0.213,0.009]$\\
@10 & Historical & 0.614 $[0.367,0.863]$ & 0.002 $[-0.014,0.018]$ & $-0.003$ $[-0.100,0.092]$\\
dyn@6 & Size-matched 95 & 1.344 $[1.121,1.556]$ & 0.098 $[0.068,0.128]$ & as @10\\
dyn@6 & Training-only 45 & 1.308 $[1.091,1.513]$ & 0.093 $[0.064,0.122]$ & as @10\\
dyn@6 & Historical & 1.261 $[1.058,1.458]$ & 0.090 $[0.063,0.118]$ & as @10\\
@7 & Size-matched 95 & 0.311 $[0.066,0.546]$ & $-0.006$ $[-0.031,0.019]$ & $-0.082$ $[-0.214,0.045]$\\
@7 & Training-only 45 & $-0.035$ $[-0.244,0.174]$ & $-0.037$ $[-0.062,-0.013]$ & $-0.110$ $[-0.234,0.014]$\\
@7 & Historical & 0.227 $[0.006,0.449]$ & $-0.016$ $[-0.041,0.010]$ & 0.003 $[-0.118,0.118]$\\
\midrule
% Deployed path (OpenVINO MiniLM + I-07b INT8 head, seam boards): analyses/I-04b_result.md, checked in I-04b_check.md.
@10 & Size-matched 95, served & 0.745 $[0.494,0.989]$ & 0.011 $[-0.005,0.027]$ & 0.011 $[-0.085,0.096]$\\
dyn@6 & Size-matched 95, served & 1.392 $[1.181,1.600]$ & 0.101 $[0.071,0.130]$ & as @10\\
@7 & Size-matched 95, served & 0.311 $[0.060,0.556]$ & $-0.004$ $[-0.029,0.021]$ & $-0.096$ $[-0.231,0.035]$\\
\bottomrule
\end{tabular}\end{center}
\textbf{Serving function.} The size-matched served rows score the boards returned by the serving function with the OpenVINO MiniLM embedder and the INT8 head; paired coverage outcomes are $59/9/3/0$ at @10 and dyn@6 and $42/20/7/2$ at @7. The OpenVINO embedder and the FP32 head reproduce the PyTorch boards on all 84 rows; INT8 changes 24 of 84 served boards (25 of 84 direct-selector boards; mean Jaccard 0.926 and 0.924, minimum 0.429) and moves one covered row at @7 only (49/71 against the PyTorch path's 50/71; a covering entry pushed to the fourth dynamic slot). The served and direct-selector boards differ on one row whose authored question flag the serving heuristic misses; on that row the served board's hedge entry ``I don't know'' matches the acceptable move ``I don't really know'', which the hedge filter does not remove on the move side, so the @10 point pass (62/71 against 61/71) rests on that row; its retention interval includes zero. On the historical artefacts the same protocol gives 61/71 at @10 and 55/71 at @7 (margins $-0.003$ $[-0.100,0.092]$ and $-0.011$ $[-0.132,0.107]$), with INT8 calibration having seen 51 of the 84 test contexts. % from analyses/I-04b_result.md and I-04_result.md; checked in I-04b_check.md and I-04_check.md
\textbf{Leaked-entry trace (@7).} Of the seven test rows the historical student covers at @7 and the size-matched 95-entry student does not, four are covered in the historical boards only by entries that entered the historical vocabulary through validation or test counts: three by ``my name is sarah'' (train/validation/test frequency 1/0/3) and one by ``i'm feeling a bit anxious.'' (2/1/1); the remaining three carry shared entries as well. Twenty-eight historical entries fall below the frequency floor on training rows alone, and nineteen of them reach it only with test rows. This trace identifies which entries supply particular metric hits on fixed boards; it is not an estimate of the retraining counterfactual, and rows 89--91 belong to one name-introduction family. % from analyses/I-05b_result.md; checked in I-05b_check.md (row ids corrected by the checker)

At @7 (continuity window) the teacher scores 4.298 distinct functions, near-duplicate rate 0.0006, same-function-pair rate 0.1735 and dynamic coverage 62/71, and the historical PyTorch student 4.095, 0.0000, 0.1893 and 56/71 with agreement 0.067. % from analyses/I-03_result.md; checked in I-03_check.md
For the three direct-selector PyTorch students, function diversity meets the 90\% threshold at @10 and dyn@6 for all three students, and dynamic coverage meets it by point estimate for none of them at @10 (size-matched 95 and historical 0.897, training-only 45 0.794) with every interval crossing zero. At @10 and dyn@6 the paired coverage outcomes (student, teacher) are $(1,1)=58$, $(0,1)=10$, $(1,0)=3$, $(0,0)=0$ for both 95-entry students and $51/17/3/0$ for the 45-entry student; every board is four pins followed by six non-core entries, so these counts are identical across the two windows by construction. The raw historical student-minus-teacher coverage difference is $-0.0986$ $[-0.1972,0.0000]$; the paired size-matched minus historical coverage difference is $0.000$ $[-0.070,0.070]$. % from analyses/I-03_result.md, I-05_result.md, I-05b_result.md; checked

\textbf{No-sibling subset (historical checkpoint).} On the retrospectively selected subset of test rows without a training or validation sibling (62 rows: 51 authored, 11 imported, no Section~\ref{sec:slate} rows), function retention is 0.936 at @7 and 0.984 at @10; coverage ratios on its 51 annotated rows are 0.867 and 0.917. Function margins are 0.155 $[-0.084,0.403]$ and 0.439 $[0.192,0.697]$; coverage margins are $-0.029$ $[-0.165,0.102]$ and 0.016 $[-0.100,0.129]$. The changed source mix prevents causal attribution to sibling overlap; intervals do not correct for retrospective selection. Of the 14 test rows with a training-or-validation context at cosine $>0.90$, 12 match training rows alone; a family-aware split is future work. % from analyses/I-09_result.md; checked in I-09_check.md (historical checkpoint only).

\textbf{Same-pool selector ablation (Table~\ref{tab:selection}; @7 rows in Table~\ref{tab:selectionfull}(b)).} Arms: relevance only (the six highest clipped cosines), MMR ($\lambda=0.6$, no answer quota or tie-break) and the full shipped policy. Paired differences over the 120 regenerated pools; counts are the paired coverage outcomes (both, first only, second only, neither). % from analyses/I-06_result.md; checked in I-06_check.md
\begin{center}\footnotesize
\setlength{\tabcolsep}{2.5pt}
\begin{tabular}{llrrrr}
\toprule
Window & Difference & Distinct fn. & Same-fn. & Coverage & Counts\\
\midrule
@7 & MMR $-$ relevance & 0.433 $[0.317,0.550]$ & $-0.035$ $[-0.046,-0.023]$ & 0.042 $[0.008,0.083]$ & 115/5/0/0\\
@7 & Full $-$ relevance & 0.467 $[0.342,0.592]$ & $-0.037$ $[-0.050,-0.023]$ & 0.025 $[-0.017,0.067]$ & 113/5/2/0\\
@7 & Full $-$ MMR & 0.033 $[-0.058,0.133]$ & $-0.003$ $[-0.012,0.006]$ & $-0.017$ $[-0.042,0.000]$ & 118/0/2/0\\
@10 & MMR $-$ relevance & 0.558 $[0.433,0.683]$ & $-0.040$ $[-0.050,-0.029]$ & 0.008 $[0.000,0.025]$ & 119/1/0/0\\
@10 & Full $-$ relevance & 0.683 $[0.567,0.800]$ & $-0.054$ $[-0.066,-0.043]$ & 0.008 $[0.000,0.025]$ & 119/1/0/0\\
@10 & Full $-$ MMR & 0.125 $[0.042,0.208]$ & $-0.015$ $[-0.021,-0.009]$ & 0 & 120/0/0/0\\
dyn@6 & MMR $-$ relevance & 0.725 $[0.575,0.883]$ & $-0.117$ $[-0.146,-0.089]$ & $-0.008$ $[-0.042,0.025]$ & 112/2/3/3\\
dyn@6 & Full $-$ relevance & 0.800 $[0.625,0.975]$ & $-0.147$ $[-0.181,-0.115]$ & 0.000 $[-0.033,0.033]$ & 113/2/2/3\\
dyn@6 & Full $-$ MMR & 0.075 $[-0.033,0.183]$ & $-0.030$ $[-0.047,-0.014]$ & 0.008 $[0.000,0.025]$ & 114/1/0/5\\
\bottomrule
\end{tabular}\end{center}
Near-duplicate pair-rate differences: MMR $-$ relevance $-0.0083$ $[-0.0119,-0.0052]$ at @7, $-0.0074$ $[-0.0107,-0.0044]$ at @10 and $-0.0222$ $[-0.0322,-0.0133]$ at dyn@6; full $-$ relevance $-0.0091$ $[-0.0127,-0.0056]$, $-0.0078$ $[-0.0113,-0.0046]$ and $-0.0228$ $[-0.0333,-0.0133]$; full $-$ MMR $-0.0008$ $[-0.0020,0.0000]$, $-0.0004$ $[-0.0022,0.0015]$ and $-0.0006$ $[-0.0056,0.0050]$. % from analyses/I-06_result.md; checked in I-06_check.md
Pool statistics: 11.64 non-core items per pool on average (minimum 6, maximum 19); 88 of 120 scenarios are in answer mode and their pools carry answer-tagged items, so ``full $-$ MMR'' compares the combined answer-quota and tie-break policy. Two pools have exactly six non-core items, so all arms select the same set there (in different orders); excluding them gives distinct functions 4.153/4.585/4.619 at @7 and 3.042/3.780/3.856 at dyn@6 (relevance/MMR/full), so the conclusion persists. Mean clipped relevance of the six chosen dynamic items is 0.354 (relevance only), 0.317 (MMR) and 0.282 (full policy): the selectors trade relevance for spread, which none of the metrics prices. The two first-seven coverage misses of the full policy (``How was your day?'' in the dinner setting) are covered at @10. The pools are verified by replay of the stored raw teacher responses, not by fresh generation. % from analyses/I-06_result.md; checked in I-06_check.md

\textbf{Candidate-pool ablation (Table~\ref{tab:pools}).} Per-arm near-duplicate rate, same-function-pair rate, mean listed hedges per board and dynamic coverage (of 71), followed by paired baseline-minus-learned differences in distinct functions and same-function pairs; positive $\Delta$DF and negative $\Delta$SFP favour the baseline. The 18-entry frequency cut-off falls inside a six-way tie at training count 5 in all three vocabularies; tied entries enter by NumPy quicksort order over vocabulary positions (2 historical, 3 training-only 45, 1 size-matched 95). The full-vocabulary arm also changes pool size. Vocabulary-order is \fn{vocab[:18]} (a pre-split-order control on the historical vocabulary; it overlaps the frequency pool on 13/18, 17/18 and 15/18 entries); random is 18 entries drawn without replacement per row (seed 0, one realisation). Intervals are uncorrected over 45 comparisons per vocabulary (12 of 45 exclude zero on the size-matched 95-entry model, 26 historical, 22 training-only 45). % from analyses/I-02_result.md and I-02b_result.md; checked in I-02_check.md and I-02b_check.md
\begin{center}\scriptsize
\setlength{\tabcolsep}{2.5pt}
\begin{tabular}{lllrrrrrr}
\toprule
Vocab. & Window & Pool & Near-dup. & Same-fn. & Hedges & Cov. & $\Delta$DF $[$CI$]$ & $\Delta$SFP $[$CI$]$\\
\midrule
Hist. & @10 & Learned & 0.0000 & 0.1619 & 0.60 & 61 & --- & ---\\
Hist. & @10 & Cosine & 0.0008 & 0.1325 & 0.15 & 65 & 0.560 $[0.345,0.774]$ & $-0.029$ $[-0.042,-0.017]$\\
Hist. & @10 & Frequency & 0.0000 & 0.1839 & 0.68 & 55 & $-0.405$ $[-0.607,-0.202]$ & 0.022 $[0.010,0.034]$\\
Hist. & @10 & Vocab.\ order & 0.0000 & 0.1701 & 0.73 & 58 & $-0.321$ $[-0.560,-0.095]$ & 0.008 $[-0.005,0.021]$\\
Hist. & @10 & Full & 0.0005 & 0.1167 & 0.13 & 64 & 1.024 $[0.810,1.238]$ & $-0.045$ $[-0.059,-0.032]$\\
Hist. & @10 & Random & 0.0021 & 0.1272 & 0.14 & 52 & 0.643 $[0.369,0.917]$ & $-0.035$ $[-0.050,-0.020]$\\
Hist. & dyn@6 & Learned & 0.0000 & 0.0937 & 0.60 & 61 & --- & ---\\
Hist. & dyn@6 & Cosine & 0.0000 & 0.0968 & 0.15 & 65 & $-0.048$ $[-0.238,0.155]$ & 0.003 $[-0.017,0.022]$\\
Hist. & dyn@6 & Frequency & 0.0000 & 0.1040 & 0.68 & 55 & $-0.107$ $[-0.286,0.083]$ & 0.010 $[-0.009,0.029]$\\
Hist. & dyn@6 & Vocab.\ order & 0.0000 & 0.0968 & 0.73 & 58 & $-0.107$ $[-0.310,0.095]$ & 0.003 $[-0.017,0.022]$\\
Hist. & dyn@6 & Full & 0.0000 & 0.0611 & 0.13 & 64 & 0.393 $[0.190,0.583]$ & $-0.033$ $[-0.052,-0.014]$\\
Hist. & dyn@6 & Random & 0.0000 & 0.0690 & 0.14 & 52 & 0.262 $[0.024,0.500]$ & $-0.025$ $[-0.045,-0.005]$\\
Hist. & @7 & Learned & 0.0000 & 0.1893 & 0.19 & 56 & --- & ---\\
Hist. & @7 & Cosine & 0.0006 & 0.1497 & 0.04 & 57 & 0.571 $[0.369,0.762]$ & $-0.040$ $[-0.060,-0.020]$\\
Hist. & @7 & Frequency & 0.0000 & 0.2166 & 0.30 & 44 & $-0.333$ $[-0.512,-0.167]$ & 0.027 $[0.011,0.045]$\\
Hist. & @7 & Vocab.\ order & 0.0000 & 0.2069 & 0.38 & 48 & $-0.238$ $[-0.429,-0.048]$ & 0.018 $[-0.000,0.035]$\\
Hist. & @7 & Full & 0.0006 & 0.1485 & 0.02 & 57 & 0.595 $[0.405,0.798]$ & $-0.041$ $[-0.062,-0.020]$\\
Hist. & @7 & Random & 0.0028 & 0.1621 & 0.06 & 41 & 0.381 $[0.167,0.595]$ & $-0.027$ $[-0.049,-0.006]$\\
\midrule
Training-only 45 & @10 & Learned & 0.0000 & 0.1714 & 0.52 & 54 & --- & ---\\
Training-only 45 & @10 & Cosine & 0.0008 & 0.1484 & 0.43 & 63 & 0.679 $[0.452,0.905]$ & $-0.023$ $[-0.035,-0.011]$\\
Training-only 45 & @10 & Frequency & 0.0000 & 0.2029 & 0.67 & 54 & $-0.702$ $[-0.893,-0.524]$ & 0.032 $[0.020,0.043]$\\
Training-only 45 & @10 & Vocab.\ order & 0.0000 & 0.1788 & 0.71 & 56 & $-0.226$ $[-0.440,-0.024]$ & 0.007 $[-0.006,0.020]$\\
Training-only 45 & @10 & Full & 0.0013 & 0.1415 & 0.39 & 61 & 0.774 $[0.560,0.988]$ & $-0.030$ $[-0.042,-0.019]$\\
Training-only 45 & @10 & Random & 0.0019 & 0.1545 & 0.43 & 57 & 0.429 $[0.190,0.667]$ & $-0.017$ $[-0.032,-0.003]$\\
Training-only 45 & dyn@6 & Learned & 0.0000 & 0.0905 & 0.52 & 54 & --- & ---\\
Training-only 45 & dyn@6 & Cosine & 0.0000 & 0.0865 & 0.43 & 63 & 0.083 $[-0.119,0.298]$ & $-0.004$ $[-0.024,0.016]$\\
Training-only 45 & dyn@6 & Frequency & 0.0000 & 0.1325 & 0.67 & 54 & $-0.571$ $[-0.738,-0.405]$ & 0.042 $[0.026,0.060]$\\
Training-only 45 & dyn@6 & Vocab.\ order & 0.0000 & 0.1103 & 0.71 & 56 & $-0.274$ $[-0.452,-0.095]$ & 0.020 $[0.001,0.039]$\\
Training-only 45 & dyn@6 & Full & 0.0000 & 0.0770 & 0.39 & 61 & 0.190 $[-0.024,0.405]$ & $-0.014$ $[-0.033,0.006]$\\
Training-only 45 & dyn@6 & Random & 0.0000 & 0.0968 & 0.43 & 57 & $-0.060$ $[-0.274,0.155]$ & 0.006 $[-0.015,0.028]$\\
Training-only 45 & @7 & Learned & 0.0000 & 0.2109 & 0.14 & 48 & --- & ---\\
Training-only 45 & @7 & Cosine & 0.0011 & 0.1791 & 0.13 & 53 & 0.512 $[0.345,0.679]$ & $-0.032$ $[-0.048,-0.016]$\\
Training-only 45 & @7 & Frequency & 0.0000 & 0.2256 & 0.36 & 43 & $-0.321$ $[-0.464,-0.179]$ & 0.015 $[-0.002,0.031]$\\
Training-only 45 & @7 & Vocab.\ order & 0.0000 & 0.2041 & 0.36 & 44 & $-0.060$ $[-0.226,0.107]$ & $-0.007$ $[-0.025,0.010]$\\
Training-only 45 & @7 & Full & 0.0011 & 0.1825 & 0.11 & 53 & 0.464 $[0.298,0.631]$ & $-0.028$ $[-0.045,-0.013]$\\
Training-only 45 & @7 & Random & 0.0017 & 0.1769 & 0.19 & 46 & 0.476 $[0.286,0.667]$ & $-0.034$ $[-0.052,-0.016]$\\
\midrule
% Clean 95 rows: analyses/I-02b_pool_v2/numbers.json (I-02b_result.md); checked in I-02b_check.md (225 intervals reproduced independently)
Size-matched 95 & @10 & Learned & 0.0000 & 0.1582 & 0.54 & 61 & --- & ---\\
Size-matched 95 & @10 & Cosine & 0.0005 & 0.1582 & 0.12 & 60 & 0.274 $[-0.012,0.548]$ & 0.000 $[-0.017,0.017]$\\
Size-matched 95 & @10 & Frequency & 0.0000 & 0.1635 & 0.58 & 56 & $-0.191$ $[-0.405,0.012]$ & 0.005 $[-0.005,0.016]$\\
Size-matched 95 & @10 & Vocab.\ order & 0.0000 & 0.1788 & 0.71 & 56 & $-0.500$ $[-0.738,-0.262]$ & 0.021 $[0.007,0.034]$\\
Size-matched 95 & @10 & Full & 0.0003 & 0.1447 & 0.17 & 57 & 0.452 $[0.179,0.726]$ & $-0.013$ $[-0.029,0.002]$\\
Size-matched 95 & @10 & Random & 0.0003 & 0.1548 & 0.24 & 49 & 0.155 $[-0.155,0.464]$ & $-0.003$ $[-0.019,0.012]$\\
Size-matched 95 & dyn@6 & Learned & 0.0000 & 0.0857 & 0.54 & 61 & --- & ---\\
Size-matched 95 & dyn@6 & Cosine & 0.0000 & 0.1135 & 0.12 & 60 & $-0.250$ $[-0.476,-0.012]$ & 0.028 $[0.004,0.052]$\\
Size-matched 95 & dyn@6 & Frequency & 0.0000 & 0.0897 & 0.58 & 56 & $-0.059$ $[-0.238,0.119]$ & 0.004 $[-0.013,0.021]$\\
Size-matched 95 & dyn@6 & Vocab.\ order & 0.0000 & 0.1103 & 0.71 & 56 & $-0.309$ $[-0.500,-0.107]$ & 0.025 $[0.005,0.044]$\\
Size-matched 95 & dyn@6 & Full & 0.0000 & 0.0992 & 0.17 & 57 & $-0.131$ $[-0.357,0.095]$ & 0.013 $[-0.008,0.035]$\\
Size-matched 95 & dyn@6 & Random & 0.0000 & 0.1032 & 0.24 & 49 & $-0.155$ $[-0.393,0.083]$ & 0.018 $[-0.005,0.039]$\\
Size-matched 95 & @7 & Learned & 0.0000 & 0.1797 & 0.19 & 50 & --- & ---\\
Size-matched 95 & @7 & Cosine & 0.0006 & 0.1746 & 0.04 & 48 & 0.155 $[-0.048,0.345]$ & $-0.005$ $[-0.023,0.013]$\\
Size-matched 95 & @7 & Frequency & 0.0000 & 0.1825 & 0.26 & 44 & $-0.167$ $[-0.357,0.012]$ & 0.003 $[-0.014,0.019]$\\
Size-matched 95 & @7 & Vocab.\ order & 0.0000 & 0.2041 & 0.36 & 44 & $-0.405$ $[-0.607,-0.202]$ & 0.024 $[0.006,0.042]$\\
Size-matched 95 & @7 & Full & 0.0006 & 0.1740 & 0.02 & 47 & 0.179 $[-0.012,0.369]$ & $-0.006$ $[-0.024,0.013]$\\
Size-matched 95 & @7 & Random & 0.0000 & 0.1820 & 0.11 & 37 & 0.119 $[-0.095,0.333]$ & 0.002 $[-0.020,0.024]$\\
\bottomrule
\end{tabular}\end{center}
Paired baseline-minus-learned coverage differences: historical cosine 0.056 $[-0.014,0.141]$ at @10/dyn@6 and 0.014 $[-0.099,0.127]$ at @7; historical full vocabulary 0.042 $[-0.042,0.127]$ and 0.014 $[-0.099,0.127]$; historical frequency $-0.085$ $[-0.183,0.014]$ and $-0.169$ $[-0.282,-0.056]$; historical random $-0.127$ $[-0.239,-0.014]$ and $-0.211$ $[-0.338,-0.085]$; training-only 45 cosine 0.127 $[0.028,0.225]$ and 0.070 $[-0.028,0.183]$; training-only 45 full vocabulary 0.099 $[0.000,0.197]$ and 0.070 $[-0.028,0.183]$; training-only 45 frequency 0.000 $[-0.056,0.056]$ and $-0.070$ $[-0.169,0.014]$; training-only 45 random 0.042 $[-0.085,0.183]$ and $-0.028$ $[-0.155,0.113]$; % from analyses/I-02_result.md; checked in I-02_check.md
size-matched cosine $-0.014$ $[-0.113,0.085]$ and $-0.028$ $[-0.155,0.099]$; size-matched full vocabulary $-0.056$ $[-0.169,0.056]$ and $-0.042$ $[-0.155,0.085]$; size-matched frequency $-0.070$ $[-0.169,0.014]$ and $-0.085$ $[-0.197,0.028]$; size-matched random $-0.169$ $[-0.282,-0.056]$ and $-0.183$ $[-0.310,-0.056]$. On the size-matched 95-entry vocabulary the frequency cut-off again falls inside a six-way tie at count 5 (one tied entry enters; the frequency and vocabulary-order pools overlap on 15/18 entries), and the learned and cosine pools share 3.8 of 18 entries on average yet reach 61 and 60 covered rows, so similar aggregate counts coexist with different pools; this comparison does not allocate coverage to selector components. The change from the training-only 45 cosine coverage advantage to an inconclusive difference on the size-matched 95-entry model compares two separately fitted checkpoints on different vocabularies; no interaction interval was computed. % from analyses/I-02b_check.md
The random pool is one realisation; its coverage deficit recurs on both 95-entry vocabularies but not on the 45-entry one. Hedge composition also differs: learned boards carry a listed hedge on 50/84 (historical), 44/84 (training-only 45) and 45/84 (size-matched 95) boards against 13/84, 36/84 and 10/84 for cosine; hedge predictions can earn coverage, and their function labels affect both diversity proxies. The hedge counts alone do not isolate a cause of the diversity differences. % from analyses/I-02_check.md and I-02b_check.md

\section{Additional observations}\label{app:additional}
% Historical timing source: docs/VERIFIED_NUMBERS.md and revision/issues_final.md I-11/I-27; no controlled cache attribution.

\begin{table}[h]
\centering\small
\setlength{\tabcolsep}{4pt}
\caption{Post-perception board-construction latency on an Apple M4 CPU (ms, pooled samples; protocol in Section~\ref{sec:latency} and Appendix~\ref{app:hw}). Jaccard is mean board-set overlap with the PyTorch path (identical boards in parentheses), not the fraction of entries replaced; p90 is descriptive of the scenario mix. Archived rows are the earlier archived-run aggregates (three repeats, 252 calls, PyTorch first) for a nominally matching historical configuration; their raw samples and original encoder weights are unavailable, so their p90 is not quoted.}
\label{tab:latency}
\footnotesize
\begin{tabular}{llrrrrr}
\toprule
Embedder & Head & p25 & p50 & p75 & p90 & Jaccard (identical)\\
\midrule
\multicolumn{7}{l}{\emph{Size-matched 95-entry configuration (\fn{b72bb3ec7254}, train-calibrated INT8 head, validated MiniLM bundle)}}\\
% from analyses/I-10_result.md (I-10_bench_v3 + v4); checked in I-10_check.md (all percentiles reproduced from raw samples)
PyTorch & PyTorch & 8.40 & 9.05 & 10.32 & 13.36 & 1.000 (84/84)\\
PyTorch & OpenVINO INT8 & 8.41 & 9.15 & 10.57 & 13.35 & 0.926 (60/84)\\
OpenVINO & OpenVINO INT8 & 3.98 & \textbf{4.50} & 5.15 & 8.52 & 0.926 (60/84)\\
\midrule
\multicolumn{7}{l}{\emph{Historical head (\fn{d6ce2e56d50b}, transductive INT8 calibration), same protocol}}\\
% from analyses/I-10_result.md section 2.5; checked in I-10_check.md (nine pooled family/backend rows match)
PyTorch & PyTorch & 8.31 & 9.22 & 11.51 & 15.82 & 1.000 (84/84)\\
PyTorch & OpenVINO INT8 & 8.34 & 9.14 & 10.75 & 13.66 & 0.913 (55/84)\\
OpenVINO & OpenVINO INT8 & 4.07 & 4.60 & 5.34 & 8.63 & 0.913 (55/84)\\
\midrule
\multicolumn{7}{l}{\emph{Earlier archived-run aggregates, historical head, archived recipe}}\\
% Archived: mac/intel/artifacts/e2e_mac.json (p50 9.6968 / 9.6508 / 5.5051); join is configuration-level only (I-04_check.md); p90 dropped per I-10_check.md
PyTorch & PyTorch & --- & 9.7 & --- & --- & 1.000\\
PyTorch & OpenVINO INT8 & --- & 9.7 & --- & --- & 0.913\\
OpenVINO & OpenVINO INT8 & --- & 5.5 & --- & --- & 0.913\\
\bottomrule
\end{tabular}
\end{table}

\textbf{Evaluation-harness timing.} The historical harness reported 10.135\,ms median (IQR 6.62--11.43; 50 scenarios: 24 classroom, 24 clinic and 2 airport, sharing ten scene strings). Its teacher timing was 7.05\,s per board (IQR 6.96--7.46; eight timing samples, all classroom scenarios), which is the generation cost the student of Section~\ref{sec:distil} removes. The unmatched workloads preclude a speed-up ratio. % issues_final.md I-27; audit_gpt.md F2; docs/VERIFIED_NUMBERS.md freeze-time rows

\textbf{Speech recognition timing.} On 2.1/4.0\,s clips, OpenVINO Whisper takes 300/403\,ms median against 252/407\,ms for \fn{faster-whisper} ($n=10$, identical transcripts).

\textbf{Exploratory Intel session.} On a Core Ultra 9 185H, head medians were 0.108\,ms for PyTorch, 0.040\,ms for OpenVINO FP32 and 0.043\,ms for INT8. In the board benchmark the MiniLM embedder always ran on the CPU and only the head was placed on the CPU (FP32) or the Arc iGPU (INT8); the PyTorch-path over OpenVINO-path board time ratios were 2.02 with the CPU head and 1.85 with the iGPU head. Each ratio changes the embedder runtime and the head together, was taken during a concurrent large download with PyTorch board timings ranging from 40 to 129\,ms across three runs, and is descriptive only. The current IR has a dynamic batch axis; a static batch-one export for the NPU has not been evaluated.

\textbf{Open-vocabulary extraction.} The context-based answer extractor activates on 2 of the 84 held-out rows and does not change the reported direct-selector criteria. % I-05_check.md; I-05b_check.md

\textbf{Other export checks.} Both training-fitted checkpoints have OpenVINO FP32 logit error below $5.8\times10^{-6}$ on all 559 contexts and reproduce all 84 PyTorch test boards. The student head uses ONNX opset~17 and OpenVINO IR; quantisation excludes the final vocabulary projection, so INT8 denotes a mixed-precision artefact. For the historical head, the maximum INT8 logit difference from FP32 is 1.05 (mean 0.23) and top-six overlap on its 360 calibration contexts is 0.950 mean and 0.667 minimum.
% Historical export checks: revision notes and mac/intel/artifacts/e2e_mac.json.
For the training-only 45-entry head, ONNX Runtime matches eager PyTorch within $7.7\times10^{-6}$ on all 559 contexts; the validation-only rule (mean top-six overlap $\geq0.8$ and minimum $\geq0.5$ against OpenVINO FP32, the repository's historical gate) rejected three of four NNCF candidates (full-INT8 validation overlap 0.786 and 0.665; head-excluded mixed preset 0.696) and kept the performance preset with the final projection excluded; candidate test logits were computed before this retrospective validation-based selection. The chosen head's test-context logit difference is at most 0.74 (mean 0.21). The ten differing test boards swap one or two dynamic entries. % from analyses/I-07_result.md; checked in I-07_check.md
For the size-matched 95-entry head the same procedure (calibration on the 391 training contexts, none from validation or test; the full-INT8 candidate's validation overlap of 0.702 fails the 0.8 gate, so the final projection is again excluded) gives test top-six overlap 0.950 mean and 0.833 minimum (59 rows at 6/6, 25 at 5/6; top-one agreement 0.976; logit difference at most 0.74, mean 0.16), 59 of 84 identical direct-selector boards with PyTorch embeddings (17 one-entry, 4 two-entry, 3 three-entry and 1 four-entry swaps; mean Jaccard 0.924, minimum 0.429), three of the membership changes involving hedge entries. Board agreement is lower than for the 45-entry head, but vocabulary content, labels and weights differ as well as width, so width is not isolated as the cause; candidate test logits were computed before the validation-only selection was applied. These agreement measurements are for the Apple M4 CPU plugin and do not transfer to Intel hardware. % from analyses/I-07b_result.md; checked in I-07b_check.md (hedge-swap and width corrections applied)
The restored MiniLM bundle has verified hashes. Archived Apple M4 CPU validation reports cosine $\geq0.9999998$ against PyTorch on 153 distinct scene/transcript texts from 84 held-out rows and all 95 vocabulary entries (maximum absolute difference $5.4\times10^{-7}$). The checker verified files and benchmark guards; numerical OpenVINO parity was reviewed from archived summaries, without retained vectors for independent recalculation. % from analyses/I-12_result.md; bounded by analyses/I-12_check.md.
The archived rows of Table~\ref{tab:latency} remain historical: bundle restoration does not reproduce their latency or establish identity with the original weights, and the 1.1\,ms gap between the archived OpenVINO median and today's session of the same recipe has no attribution. % from analyses/I-10_result.md; checked in I-10_check.md
\textbf{Benchmark caveats.} The p90 values (8.52\,ms OpenVINO, 13.4\,ms for both PyTorch-embedder paths) describe this scenario mix and cache/load protocol; a few long-context question rows contribute disproportionately. % from analyses/I-10_result.md; checked in I-10_check.md
The earlier archived benchmark of the historical head (PyTorch first, three repeats) reported 5.5 and 9.7\,ms medians for a nominally matching configuration; its larger PyTorch-path tail recurs under that recipe today, but the counterbalanced PyTorch and INT8-head tails nearly coincide (13.36 against 13.35\,ms), so the tail is protocol-sensitive and not attributed to either head. % from analyses/I-10_result.md and I-04_check.md; checked in I-10_check.md (mechanism-neutral wording required)
One session of the archived recipe on the historical artefacts today gives medians of 9.7 (PyTorch), 9.3 (INT8 head, PyTorch embedder) and 4.4\,ms (OpenVINO) with a PyTorch-path p90 of 21.9\,ms against 12.3\,ms for the INT8 head, reproducing the archived pattern in one session; the cause of the archived asymmetry is not established, since order, repeat count, eviction, reset and static warm-up change together between the recipes. Five of the eighteen size-matched-configuration sessions were flagged for another process at or above 40\% CPU in a bracketing snapshot; restricting to unflagged sessions moves the size-matched medians by at most 0.05\,ms and the p90 by at most 0.15\,ms, but snapshots cannot see bursts inside a session. Per-position medians differ by at most 0.62\,ms with no consistent direction, without establishing elimination of order effects; the single isolated process per backend ($n{=}1$) is about 0.3--0.4\,ms slower at the median for the PyTorch-embedder paths and not separable from drift. The first timed pass through a freshly loaded PyTorch embedder has a round-one p90 of 21.8--25.1\,ms against 12.6--13.4\,ms afterwards, with no comparable tail pattern for the OpenVINO embedder; the mechanism is not identified. For the OpenVINO path only 18 of 84 scenarios contribute any sample above the pooled p90, and the six slowest are all question rows on which the open-answer extractor encodes additional spans. The size-matched-configuration timing and quality records share their artefact and source hashes and their first-round boards, but the current working-tree \fn{infer.py} differs from the recorded one, so the join describes the recorded pipeline rather than the current tree; later-round boards were not archived. Section~\ref{sec:distil} and Appendix~\ref{app:uncertainty} report exported-path quality for the size-matched 95-entry model and historical artefacts, and the head sizes cover the head IR only, not the full system footprint. % from analyses/I-10_result.md; checked in I-10_check.md (caveats and infer.py drift)

\section{Use of large language models}
LLM coding agents assisted with code, tests, analysis scripts and notes under author direction, and LLM tools assisted with manuscript drafting and editing. In revision, LLM agents ran the additional analyses reported here (window re-scoring, training-only and size-matched retraining, candidate-filter and same-pool comparisons, requantisation, export validation, serving-path scoring and the latency benchmark), and a second agent checked the reported analyses within the reproduction scope documented in its checker report; provenance for each number is recorded with the artefacts (Appendix~\ref{app:repro}). The system itself uses a local language-model teacher and a vision--language model. The authors are responsible for the implementation, measurements, citations and final claims. % I-22

\end{document}
